% Canonical Life Simulation pre-experiment method manuscript.
\documentclass[10pt]{article}

\usepackage{emergentwisdom-longform}
\usepackage{mathtools}
\usepackage{booktabs,tabularx,array,float}
\usepackage{enumitem}
\usepackage{xcolor}
\usepackage{graphicx}
\usepackage{tikz}
\usetikzlibrary{arrows.meta,positioning,fit}
\usepackage{hyperref}
\usepackage{multicol}
\geometry{bottom=0.7in}

\newcolumntype{Y}{>{\raggedright\arraybackslash}X}

\hypersetup{
  hidelinks,
  pdftitle={Life Simulation: Learning from Worlds and Their Construction},
  pdfauthor={Henrik Westerberg},
  pdfsubject={A framework for explicit time-indexed world processes, candidate dynamics, joint process prediction, and prospective learning},
  pdfkeywords={life simulation, world models, hybrid processes, process supervision, generative inversion, prospective learning}
}

\newcommand{\system}{\textsc{LifeSim}}
\newcommand{\MeaningModel}{\textit{Meaning Model}}
\newcommand{\Book}{\textit{The Book of Conditions}}

\ewlongtitle{Life Simulation}{Learning from Worlds and Their Construction}
\author{Henrik Westerberg\\
  \small{henrik.westerberg@emergentwisdom.org}}
\ewpublicationdate{September 5, 2026}

\begin{document}
\raggedbottom
\maketitle

\begin{abstract}
Life Simulation proposes learning to maintain continuing worlds, understand
their processes, and improve the representations and capabilities used to act
within them. Constructed worlds become rich linked learning material: one
history supplies observations, process paths, concepts, alternatives, and
changing interpretations. The Meaning Model supplies their shared semantic
interface and preservation contracts; Life Simulation advances those worlds
and turns their histories into learning objectives.

A proposed process sensorium predicts native observations and explicit process
state together, learning persistence, accumulation, delay, and change across
resolutions. Forward construction supplies controlled worlds and varied
renderings for learning the inverse: reconstruct only what observations
support. Applied to licensed sources, this could progressively build a
connected corpus model with shared identities and retained disagreements.
World-history supervision teaches what happens; construction-history
supervision teaches which questions, distinctions, and revisions improve an
account, including how to deepen or reorganize a personality model.

The learner's organization can change too. Fractal Intelligence retains and
reorganizes reusable capabilities; Alien-world search revises the categories
that direct mechanism exploration. Their consequences can train subsequent
choices. In the recommended combined training, a continuing Reader develops
an Understanding Graph under Entangled Alignment's evaluative orientation.
Native spans, world-grounded visuals, process histories, and Reader corrections
jointly connect understanding people to decisions. Exact concept versions
identify declared meanings; information cutoffs delimit available evidence.
Later evidence supplies correction targets without entering earlier inputs.
Durable care remains an objective, not an automatic consequence.

A process path may be continuous, discrete-event, piecewise, stochastic, or
hybrid, advanced by language models, simulators, estimated functions, or their
compositions. The path is required; a compact governing law is an optional
hypothesis. Narrative is the first controlled laboratory, not the scope of the
proposal. This paper develops the method and specifies experiments for
evaluating its learning hypotheses.
\end{abstract}

\keywords{Life Simulation, Time-Indexed World Models, Hybrid Processes,
Process Sensorium, Generative Inversion, Prospective Learning}

\ewfrontmatterend

\section{Introduction}
\label{sec:introduction}

Text, measurements, and actions arrive as fragments. The systems that produced
them continue between observations. A person changes while no message is
written; an institution accumulates commitments between announcements; a
scientific field reorganizes before a new vocabulary becomes public. A model
may infer this persistence from the visible stream. Life Simulation makes it
an explicit learning target: maintain an account of which processes exist,
how they changed, what was known at a cutoff, and which later observation
corrected the estimate. It also makes the choice of that account learnable:
which questions to ask, which detail to open, and when a representation or
capability organization needs revision.

Life Simulation makes a selected set of those processes explicit. At minimum,
a simulation contains a world state, time-indexed process paths within that
world, a mechanism that proposes or computes their continuation, and an
observation mechanism that exposes only part of the state. Natural language is
one observation stream among others. A continuous function can be useful, but
it is not the definition: births, discoveries, messages, commitments, and
failures are discrete; fatigue and cash may vary continuously; many useful
models combine both.

There are two linked ambitions. One is to make a world useful while it is still
developing: its processes continue, actions have consequences, and detail is
opened where an interaction needs it. Such a world can support play, creative
work, or exploration without first becoming a training corpus. The other is to
learn from the histories those worlds expose. A learner would not only predict
the next observation, but maintain an account of what produced it, anticipate
the effects of a change, and revise that account when evidence disagrees.
Neither use requires an exhaustive world. Their value is tested separately:
useful interactive continuity does not establish a learning advantage, and a
learning result does not by itself establish a useful simulation.

For a concrete example, suppose a story says that a promised payment has not
arrived. A world model can retain the commitment, deadline, nonreceipt, and
affected people. A continuation follows what changes when the delay persists
or payment arrives. A timeline or filmed scene can expose selected
consequences, while a Reader's linked interpretation asks whose plans depend
on the money and revises its answer as evidence appears. Those aligned records
can then become learning material: the text, the visible events, the processes
behind them, and the developing interpretation. An author can create such a
history; reconstructing it from an existing text must leave unsupported
details unresolved. This is the construction-to-learning chain developed
throughout the paper.

The alignment-facing ambition gives this combination a particular purpose:
learn to understand people well enough to treat them well, and learn to make
that understanding matter in decisions. Recognizing an emotion is not enough.
An agent needs an account of what the person wants, believes, has experienced,
and can change, together with how its own intervention may affect that life.
The proposed solution joins that contextual understanding with the formative
evaluative learning of Entangled Alignment, rather than assuming that accurate
prediction will produce care by itself. Life Simulation is the integration
method for this proposal: it connects the program's representations, reading,
problem-solving, and training methods through common evolving histories.
It does not replace the companion papers or claim that their combination has
already solved alignment.

This paper separates two responsibilities that operate together. The
Meaning Model asks how an explicit world can be represented, progressively
opened, interpreted, and rendered without losing identity or confusing one
perspective process's account with world authority~\cite{westerbergMeaning2026}. Life
Simulation uses that shared representation: typed world facts and quantities,
Cuts and abstract concept models, Event descriptions, story-passage text, and
externalized understanding have exact addresses and links while retaining
different authority. They can be constructed and revised together; declared
dependencies identify which accounts and renderings need reconsideration when
a world record changes. This construction contract belongs to Meaning Model.
Life Simulation asks what can be done with the evolving histories: how to advance
them, search for dynamics, infer hidden paths from observations, train models
on synchronized streams, and correct them using later evidence. A proposed
continuation or new process distinction remains subject to Meaning Model
acceptance and revision rules; this is not a handoff from a finished world.
The Meaning Model is the selected interface here, but the temporal and learning
tests can also be applied to another world representation that supplies equivalent
identity, time, provenance, and uncertainty contracts.

The central hypothesis is not that every aspect of reality follows one
equation. It is that useful regularities become easier to discover and learn
when observations are accompanied by explicit process histories at the right
resolution. A recurring pattern is slow trajectory, faster perturbation, and
adaptation. This rhythm may appear within a conversation, a life, an
organization, a technological lineage, or an ecosystem. Jump diffusion is one
candidate family for a continuous state with exogenous shocks; it is not a
universal law. Planned transitions, clustered events, absorbing states, and
adaptations that never settle require other families.

The proposed bootstrap has seven stages. They are ordered so that hidden-state
recovery is tested where the hidden state is known before the same estimator is
trusted on real chronology. There are two learning stages: controlled worlds
first support an inverse estimator; its source-grounded reconstructions then
contribute to the larger multimodal and Reader-training proposal.

\begin{enumerate}
  \item \textbf{Construct forward.} Build or import an explicit world at a
    declared resolution and retain the status and provenance of each selected
    process. Progressively refined books and game worlds provide candidate
    histories whose observations can be rendered in several forms.
  \item \textbf{Vary the dynamics.} Advance matched worlds under direct model
    judgment, candidate functions, parameter changes, or controlled random
    inputs, preserving which mechanism produced each path.
  \item \textbf{Learn the inverse.} Hide simulator-privileged state and infer a
    distribution over process paths and later outcomes from observation
    prefixes.
  \item \textbf{Project onto licensed material.} Apply the estimator to fiction,
    biography, dialogue, historical records, or other chronological streams
    without turning attribution into observation. At corpus scale, a continuing
    EA Reader can develop an Understanding Graph across these source-grounded
    models rather than attach isolated reactions to unrelated passages.
  \item \textbf{Compile synchronized targets.} Use \emph{Imagining the Corpus}
    to align exact native spans with model-grounded images or video, selected
    Meaning Model process records, and Reader interpretations for joint
    prediction~\cite{westerbergImagining2026}. Section~\ref{sec:alignment}
    develops this construction-to-training chain.
  \item \textbf{Test prospectively.} Freeze each information cutoff and forecast,
    then score it only after later reports, actions, measurements, or
    interventions arrive.
  \item \textbf{Transfer the discipline.} Test the same endpoint-versus-history
    distinction on Fractal Intelligence traces, interactive worlds, and alien
    mechanism searches rather than treating narrative as the final domain.
\end{enumerate}

Narrative is the first controlled laboratory because an authored world can
provide a complete accepted history, private beliefs, counterfactual branches,
and multiple renderings without pretending that an inference about a real
person is ground truth. Narrative is the first construction laboratory; the
reservoir in Section~\ref{sec:first-experiment} is the first proposed controlled
learning experiment, isolating a smaller claim before the richer setting.
The companion Meaning Model paper specifies the new
construction of \Book{}, its world-building protocol, and its completed
manuscript and partial construction witness. Independent read-back and the
matched story comparison remain prospective. Life Simulation can use released
histories as experimental material or supply continuations during a new
construction. The companion still governs how those proposals are represented,
accepted, revised, and rendered.

Table~\ref{tab:program-roles} identifies the nine companion contributions and
the role Life Simulation plays in connecting them. The combination is a
research architecture, not nine finished modules that every experiment must
instantiate.

\begin{table}[!htbp]
\centering
\small
\begin{tabularx}{\textwidth}{@{}p{0.24\textwidth}Y@{}}
\toprule
Component & Role in the proposed program \\
\midrule
Meaning Model & Constructs and progressively refines explicit worlds, concepts, perspectives, and narrative surfaces \\
Life Simulation & Evolves time-indexed processes, tests candidate dynamics, and learns from synchronized process and observation histories \\
Sema & Identifies exact definition and grounding versions used across world, interpretation, and solver records \\
Understanding Graph & Retains a continuing Reader's questions, interpretations, disagreements, and revisions with source links \\
Entangled Alignment & Specifies formative Reader training with a recurrent evaluative orientation and tests whether it shapes durable conduct \\
Imagining the Corpus & Pairs exact native spans with evolving explanatory visual or schematic tracks; provides the synchronization pattern for the combined corpus \\
Temporal Hindsight Learning & Enforces cutoff blindness and turns later evidence into labels for earlier predictions \\
Fractal Intelligence & Searches and revises conceptual decompositions; retains evaluated capabilities for reuse and learning from execution \\
Ontology of the Alien & Connects causal-world diversity with a revisable candidate ontology that directs further mechanism search \\
Substrate Language Modeling & Tests the distinct architectural choice of requiring lexical prediction to pass through a predicted substrate \\
\bottomrule
\end{tabularx}
\caption{Life Simulation and its nine companion contributions. Shared histories
connect their roles; successful integration remains to be demonstrated.}
\label{tab:program-roles}
\end{table}

The contribution is a proposed change in what the learner must maintain and
improve. A process sensorium jointly predicts observations and the continuing
processes they expose. Two histories give that learner different tasks:
represented-world history teaches what happens, while model-construction
history teaches how to choose, open, and revise an account. A planned ending
may guide a writer's construction but cannot enter a character's earlier
forecast. Paired histories could teach when to change the questions used to
understand a person, not only the values in a fixed personality profile.
Meaning Model preserves the commitments across these changes; Life Simulation
defines the learning views and their cutoffs.

This is also a data-construction contribution. One narrative can support linked
examples across people, concepts, processes, timescales, observation cutoffs,
and alternative continuations. Constructed worlds supply authored state and
varied renderings for inverse learning; corpus interpretation supplies
evidence-constrained estimates and connects supported identities across
sources. Descriptions, prose, visuals, and understanding histories can then
train together. Progressive detail expands this material where useful, while
corrections can change the connected corpus model itself. This could produce
a very large training corpus whose learning targets include both temporal
behavior and the development of its interpretation. Generated variants remain
correlated, inferred values remain uncertain, and invented detail is not
historical evidence. Coverage, cost, and transfer must be assessed alongside
volume.

A further hypothesis makes capability organization part of this changing
world. The Meaning Model can provide the evolving world, visible evidence,
available options, actions, consequences, and perspective-specific readings.
Fractal Intelligence can provide instrumented Events tied to persistent,
concept-defined callable Solvers. Proposed decompositions are assessed for
necessity, independence, universality, and completeness within their declared
scope; execution can retain
a reusable capability or motivate splitting, merging, or reparenting it.
Learning therefore concerns which capabilities to form and compose, as well
as which route to take. This requires that telemetry is causally tied to
execution rather than reconstructed afterward; Section~\ref{sec:problem-solving-events} specifies
the proposed adapter and controls. Alien-world search changes another object:
the categories that direct exploration. A diagnosis of repeated mechanisms can
request a different causal relation, whose simulated consequences inform the
next search and category revision. These are proposals for learning how to
organize competence and exploration, not only retaining their trajectories.

In the recommended combined training, a continuing Reader's evaluative
orientation is central to what this growing competence is for. Its
Understanding Graph records how evidence changes an interpretation of people
and consequences; synchronized native, visual, and process targets connect
that interpretation to the world. Reader updates may guide subsequent
intervals without rewriting earlier sources or forecasts. Exact grounding
versions identify the concepts used, and later evidence can label earlier
judgments while remaining outside their inputs. The operational coupling is
the passage from an accepted world through rendering, inference,
interpretation, action, and correction into subsequent learning. Whether it
improves prediction, representation choice, reusable competence, or conduct
requires separate matched tests; none is established by connecting the records.

Section~\ref{sec:temporal-object} states
the learning loop alongside the world it concerns; Section~\ref{sec:evaluation}
then gives a small, implementable first experiment. Candidate dynamics,
adaptive resolution, narrative, execution histories, and human-facing learning
extend that experiment rather than becoming prerequisites for running it.

\section{Explicit Worlds with Time-Indexed Processes}
\label{sec:temporal-object}

The following object separates a world's evolution from what a learner can
observe. Its projections then supply the joint learning task, before the
paper considers particular dynamics or applications.

\subsection{The simulation object}

Let $W_t$ denote the accepted explicit world at time $t$. For interval
$I=(t,t+\Delta]$, let $\widetilde X_I$ be a proposed path for selected
continuous, allocated, or derived process state; let $\widetilde E_I$ be the
proposed discrete Events; and let $Y_I$ be the observations made available from
the interval. A Life Simulation supplies a transition mechanism, a
world-interface update, a selected projection, and an observation mechanism:

\begin{align}
  (\widetilde X_I,\widetilde E_I)
    &\sim G_{\phi}(\,W_t,X_{\leq t},E_{\leq t},u_I,Q\,),
    \label{eq:life-transition}\\
  (W_I,W_{t+\Delta})
    &=U_{\mathcal V^{(k)}}(W_t;\widetilde X_I,\widetilde E_I),
    \label{eq:world-interface-update}\\
  E_I&=\mathcal E(W_I),
    \label{eq:event-projection}\\
  X_I&=\Pi_Q(W_I,E_I),
    \label{eq:process-projection}\\
  Y_I
    &\sim O_{\psi}(W_I,X_I,E_I;h),
    \label{eq:life-observation}
\end{align}
Here $G_{\phi}$ proposes an entire interval path and its discrete Events, not
only an endpoint. $U_{\mathcal V^{(k)}}$ is the selected world interface's
validated update under registry version $\mathcal V^{(k)}$; it returns the
accepted interval history $W_I$ and successor world $W_{t+\Delta}$. The Event
path $E_I=\mathcal E(W_I)$ is an index or projection of Event records already
in that accepted history, not independent state or a second authority. In the
Meaning Model instantiation, its own construction and acceptance contract
supplies this update. $Q$ declares which process
identifiers, fields, structural perspective roots, and temporal resolution are selected, and
$\Pi_Q$ projects their accepted path as $X_I$. Thus $X$ is an attachment to or
view of accepted world history, not a second authority that can contradict
$W$. $O_{\psi}$ receives only the now-defined accepted interval, selected
projection, Events, and access policy $h$. The external inputs are
$u_I$. Neither $G_{\phi}$ nor $O_{\psi}$ must be closed form: each may be a
language-model proposal, deterministic code, a stochastic process, an
empirical table, a learned model, or a pipeline of several mechanisms.

Equation~\ref{eq:life-transition} describes a generated continuation, but not
every accepted path is generated by $G_{\phi}$. A recorded trajectory may enter
the same update through an observation adapter, retaining its observed status
and never being redescribed as a model output. Before either route is used, a
process declaration fixes its world address, temporal support, value domain and
unit when one exists, and sampling and missingness policy. Its uncertainty is
also factored rather than pooled: transition stochasticity, uncertainty about
latent state, parameter uncertainty, measurement error, population
heterogeneity, and model discrepancy answer different questions and remain
separately reportable.

The indispensable object is the versioned path, not the equation used to
produce it. Continuous coordinates retain units and sampling rules. Discrete
events retain intervals, participants, provenance, and state changes. Hybrid
processes relate the two: a discrete injury changes a continuous capacity path;
a threshold crossing creates an event; an institutional decision changes the
regime governing a financial stock. Unknown, unmodeled, and unresolved states
remain legal rather than being filled with invented precision.

The native stream $Y$ is not subordinate to the explicit track. It is the text,
image, audio, action, sensor reading, transaction, or document a system would
actually observe. The process projection $X$ and discrete Event path $E$ are
selected parts of the accepted history whose cost and value must be measured. A latent model can be a
stronger comparator; explicit state earns its place only when it improves
prediction, transfer, control, correction, or audit at matched total cost.

\subsection{Process fields, allocation, and events}

Three numerical and structural roles must remain distinct. A state coordinate
is a numeric process field with an independently declared unit or measurement
protocol; other process fields may be categorical. Concurrent fields do not
sum merely because they describe the same Thing. An allocation divides one
named unit among mutually exclusive
answers and therefore sums to one. Events occur within or across processes and
may change several coordinates at once; the set of events is not normalized.
When the Meaning Model is used, its Cut record supplies the allocation case and
typed world data supplies the coordinate case. Life Simulation does not add a
third kind of semantic number.

This distinction permits both literal conservation and qualitative modeling.
Cash, elapsed time, energy, or population may obey accounting identities in
their native units. Attention or attributed causal responsibility may be
represented as local shares under a declared question. Health and standing are
concurrent categorical or externally measured process fields, not numeric
coordinates unless an independent unit or protocol is supplied and not
competing pieces of a person. A discovery can
perturb science, industry, and politics simultaneously without forcing their
effects into a simplex; a separate impact allocation is possible only after a
named perspective states what total unit is being divided.

\subsection{Slow trajectories, shocks, and adaptation}

Processes are nested by temporal resolution. Slow paths summarize persistence
or drift; faster events perturb them; still finer events may be linked to the
local shape of one perturbation. A perspective-specific causal explanation remains a
separate attribution. A long process can be opened into disjoint periods
for aggregation and also linked to overlapping episodes such as illness,
partnership, or war. The periods are an accounting partition. The episodes are
causally or narratively coherent arcs and need not tile the parent.

A shock is defined relative to a selected process and resolution: an event or
short interval that produces a material departure from its recent path or
expected continuation. Its reach is the set of processes it measurably or
interpretively changes. Anticipation belongs to a named perspective and cutoff. Before the outcome, a
direction Cut declares whether its weights are a proposal policy or forecast
probabilities over exclusive continuations at a stated horizon. Proposal
weights guide construction; forecast probabilities express predicted
likelihood and can be scored once the outcome is observed, even before
calibration has been established. Calibration is evaluated empirically across
forecasts.
Adaptation is the subsequent path. Useful descriptive shapes
include drift, recovery, consolidation into a new baseline, instability or
cascade, terminal transition, and anticipated transition in which adjustment
begins before the event.

Operational descriptors make those shapes comparable without declaring a
governing equation. Shock magnitude is only a change in a native-unit measured
field, total variation between compatible Cuts, or a named categorical
transition---never a freely chosen scalar. Other descriptors include recovery
half-life; cascade fan-out; delay from shock to slow-state change;
and dependence of later event hazard on the preceding state. Every descriptor
retains its process, resolution, sampling rule, evidence status, and uncertainty.

One candidate multiscale pattern gives fast events durable influence. Repeated
events write to a memory or resource stock; the stock crosses a threshold; a
slower process field or regime changes; and that change feeds back into later
sensitivity, recovery, available actions, or event hazard. Fatigue, trust,
institutional backlog, and accumulated anomalies can all exhibit versions of
this shape, but they need not use the same representation. Fatigue or backlog
may be categorical or independently measured stocks; accumulated anomalies may
be counts; attributed trust is an Event-and-Cut trajectory contained within the
attributing actor's or observer's perspective process unless an independent
measurement protocol is supplied. The shared diagram is not
evidence for a shared law. Each instance
must beat no-memory, linear-accumulation, and direct-model baselines on held-out
paths before the feedback pattern is retained.

The same vocabulary can be applied recursively, but not mechanically. A war
may be a society-scale shock, a years-long episode in one person's life, and a
sequence of minute-scale shocks in a single encounter. Whether these views
compose is a declared modeling question. The method does not infer an
individual's trauma from a collective label, does not invent an ancestral cause
to close a remainder, and does not treat a plausible jump-and-recovery story as
a validated psychological law.

\subsection{Local models and perspective}

When agents act from incomplete information, the relevant state is not only
the world history but the history available to that actor. Let
$\mathcal M_a(t)$ denote the projection derived from representational Events
under actor $a$'s governing perspective root at $t$, together with records
permitted by the access policy,
and let $\mathcal M_a^*(t)$ denote the unknown or authored operative organization
actually used in a decision. An observer $o$ may maintain one or more estimates
$\widehat{\mathcal M}^{*}_{o\rightarrow a,j}(t)$. World state, actor projection,
operative model, self-report, and observer hypotheses are separate objects. A
simulated decision for $a$ receives only records accessible through
$\mathcal M_a^*(t)$ and actually observed or delivered by the cutoff. Mere
co-presence or reference does not reveal every Event field or another actor's
inner state; disagreement with the world is misperception, while
disagreement between perspective projections is a relation to be modeled rather than averaged
away.

For an actual person, four epistemic objects must not collapse: external
evidence about conduct and circumstance; several observers' competing
hypotheses about the person's operative organization; the person's dated
self-report; and whatever operative organization remains unknown. An
Understanding Graph trace externalizes a selected set of questions,
interpretations, reports, and revisions. It is not asserted to be a complete
mind, and silence in that trace is missing cognitive annotation rather than the
absence of thought.

This makes theory of mind temporal. The target is not a single trait label but
a calibrated path over beliefs, wants, appraisals, opportunities, and estimates
of other actors. Several paths may remain compatible with the same behavior.
They earn a role only if they improve later prediction or intervention relative
to summaries, static traits, and an equally informed latent baseline.

\subsection{A proposed process sensorium}

The larger learning claim is to predict native observations and selected
explicit process state together. Given cutoff history $H_t$, a student may be
trained toward

\begin{equation}
 p_{\theta}\!\left(
   Y_{(t,t+\Delta]},X_{(t,t+\Delta]},E_{(t,t+\Delta]},\Delta\mathcal V
   \mid H_t,Q,\mathcal V^{(k)}
 \right),
 \label{eq:joint-process-prediction}
\end{equation}

where $H_t$ is the exact cutoff-visible native and accepted-world history;
$Y$ is the native text, image, audio, measurement, or action stream; $X$ is the
selected coordinate, allocation, and derived-process path; and $E$ is the
selected index of accepted discrete Events in that history. The query $Q$ fixes the process
identifiers, fields, applicable root paths, and resolution to be predicted, as in
Section~\ref{sec:temporal-object}. The registry version $\mathcal V^{(k)}$ is
frozen at the cutoff, and $\Delta\mathcal V$ contains proposed model revisions.
New or revised concept definitions, process lenses, world-data schemas, Cut or
Binding profiles, and grounding versions require a separate versioned model
revision before later predictions may use them. Ordinary Event, Cut, and Binding
instances using the frozen registry enter through validated world updates;
they do not change the registry. Unknown, unmodeled, and uncertain are valid
outputs; the model is not rewarded for inventing precision.

The explicit process track is a proposed sensorium: training on emotional,
spatial, social, institutional, or problem-solving trajectories asks a model to
maintain estimates of persistence, accumulation, threshold, delay, recovery,
coupling, and regime change instead of mentioning a state only when a nearby
token makes it salient. This is a testable capability claim, not a claim of
consciousness or direct perception. Simulator-privileged targets during
training, posterior inference from ordinary observations, and explicit sensors
at deployment are three different conditions and must be reported separately.

The intended numerical intuition is an ability to connect visible situations
to these structured quantities without reconstructing the entire account from
scratch on every query. It is not access to a unique set of numbers hidden
behind every human life. A pound balance retains its external unit; an emotion
Cut retains its question, sibling-relative shares, remainder, and perspective.
Their changes can support anticipation and correction without making the
numbers interchangeable. The proposed new sensory faculty is functional:
notice accumulation, an approaching limit, or a changing balance of concerns,
and produce a calibrated estimate when asked. Unfamiliar scales and explicit
numerical prediction tests must distinguish this capacity from fluent talk
about change.

The aligned track is compared with the five-level frontier and the
time-permuted and state-scrambled controls in Section~\ref{sec:evaluation}.
A benefit shared with scrambled tracks does not establish a synchronization
advantage. The additional gain of the aligned track over its matched scrambled
control is the relevant test; both effects may coexist.

If the capability transfers across unrelated domains, operates on each
process's native time and units, and survives those controls, it may be useful
to call it \emph{process-native intelligence}: competence at maintaining and
reasoning over evolving processes rather than only describing their latest
surface. The name is conditional on the evidence and is not a claim of general
intelligence.

Persistent residual error may motivate a proposed Event boundary, process
lens, measured field, Concept, decomposition, or grounding. Such growth earns acceptance only if it
improves recomposition, held-out prediction, intervention, or transfer enough
to pay its acquisition and maintenance cost. No present artifact trains or
validates this faculty.

\subsection{Synthetic bootstrap and generative inversion}

Synthetic semantic worlds provide the controlled first rung because the
generator knows their accepted Events, structurally separated actor
perspectives, actions, rejected siblings,
and later consequences. Independently varied renderers can express one world
as prose, dialogue, images, audio, or game traces. A state-blind inverse model
then reconstructs only what a prefix supports. Cross-renderer recovery tests
whether it learned represented process structure rather than one writer's
phrasing.

The learning route is

\begin{equation}
 \begin{aligned}
  \text{accepted structured world}
  &\rightarrow\text{varied observable renderings},\\
  \text{rendering prefix}
  &\rightarrow\text{distribution over world records}
  \rightarrow\text{predicted later behavior}.
 \end{aligned}
 \label{eq:generative-inversion-method}
\end{equation}

This is analysis by synthesis. Narrative quality is part of the laboratory:
worlds whose renderings are judged coherent and meaningfully surprising have
passed a generative-coherence test. They have not established that their
authored psychology is true. After identifiability is tested inside synthetic
worlds, the representation can be tested on licensed fiction, biography,
dialogue, historical records, and consented longitudinal data with uncertainty.
Later actions, reports, measurements, and interventions score earlier
distributions. The compact rule is synthetic-centered bootstrapping and
reality-centered validation.

Writing from a model and reconstructing a model from a book are related but
not interchangeable achievements. An author may choose an unspecified motive
or invent the missing event; an inverse reader must retain alternatives where
the text does not decide between them. An unreliable narrator or inconsistent
account may require several perspective-local reconstructions rather than one
complete world. Neither direction is declared universally harder. Forward
success supplies controlled examples for inverse learning, not proof of its
success on arbitrary books. The eventual corpus-wide ambition is repeated,
source-grounded reconstruction with selective human correction, uncertainty,
and a recorded cost, not exhaustive recovery of every author's hidden model.

\paragraph{A connected Meaning Model of the corpus.}
The larger ambition is a progressively constructed Meaning Model of what the
accessible corpus describes, not only an independent reconstruction for each
book. Accounts of the actual world can connect to shared people, places,
institutions, and Events where evidence supports the identification. A
biography, a letter, and an economic history may refine different parts of the
same period. Each contribution retains its source, access conditions, and
uncertainty; incompatible accounts remain contestable rather than being
merged into an apparently certain history. Fictional and counterfactual
worlds remain separate, even when they refer to real places or historical
people.

Coverage can grow across world history and geography while resolution deepens
selectively. A coarse account of a city over a century may later open into
decades, institutions, relationships, and individual Events, with links to
processes elsewhere. New sources can also revise the coarse account rather
than merely fill it in. A newly read letter may change the model's account of
an earlier year; that revision must not silently enlarge the information
available to an actor at that earlier time. Updating the external Meaning Model is distinct from
training model weights: selected snapshots and construction histories can
subsequently supply training examples under their own cutoffs. Whole-corpus
coverage is the direction of growth, not a claim that missing history or
private experience becomes knowable.

Authored worlds and actual-world reconstructions therefore serve different
training purposes. Fiction permits chosen events, explicit inner lives, and
controlled alternatives; acceptance establishes fictional canon. An account
of reality must answer to independent evidence, preserving unknowns and
attributed interpretations instead of inventing facts to complete the graph.
Both regimes can supply construction histories and represented-world
histories, the separate learning views in Section~\ref{sec:learning}.
Synthetic examples can teach representation and simulation skills, but their
consistency cannot validate claims about reality.

Generation also provides a way to explore candidate mechanisms. A proposed
process must sustain a history whose consequences can be inspected, rather
than merely sound plausible in a description. Failure may motivate a different
transition rule or a revision of the represented processes and concepts, as
allowed by the sensorium's model-revision output. New histories then test the
revised account; selected correction histories can later train a successor
learner. This closes the loop from constructing worlds to improving how they
are understood and constructed. Coherence can select interesting candidates,
but independent observations and interventions must decide whether any
candidate explains reality. The loop is a research proposal, not a general
algorithm for discovering laws.

\section{The Meaning Model as World Interface}
\label{sec:meaning-interface}

This paper assumes, rather than re-derives, the Meaning Model
\cite{westerbergMeaning2026}. Its stable interface consists of registered
Concepts and Things, time-bearing Events, typed Bindings, local Cuts, and
Realizations, each with provenance. Every modeled Thing is associated with a
lifecycle Event. Events can carry typed world data, participate in longer
processes, and update selected coordinates. A Cut divides one declared unit
among mutually exclusive children and a remainder. Every Cut has an Event
parent and inherits context through edges declared \textsf{authority-parent},
not through every link to a parent. These paths must resolve to the same
nearest context root; otherwise the record is rejected as ambiguous. Temporal
containment, occurrence, reference, evidence, and grounding alone confer
neither authority nor access. A Realization links a fragment to a versioned Concept grounding. These
distinctions let Life Simulation ask what changes without redefining what the
records mean.

The selected roots make that ancestry explicit. Accepted world Events descend
from the Universe's lifecycle Event, History. That temporal ancestry does not
override a nearer actor-local authority root. A belief, appraisal, forecast, or
estimate attributed to an actor is a distinct Event in a named inner Event
process beneath that actor's lifecycle. Understanding Nodes descend
from a named understanding-process root, and Document Nodes descend from a
named document root. A Document Node with role \textsf{story-passage} is the
addressable unit used by rendering and read-back; it is not an additional
record form. An Event may carry an optional readable description field for that
exact Event version. This field is neither an Understanding Node nor a Document
Node and has no independent authority.

Every record query returns its permitted parents and authority-parent paths to
the governing context; temporal ancestry is identified separately. Access and
the evidence cutoff are checked before traversal or serialization. A denied
reference exposes neither private content nor private ancestry. A query or
export containing a Cut or component returns the complete permitted local
Cut---parent Event, question, unit, keyed sibling answers and weights,
remainder, optional conditioning address, constraints, and recomposition
rule---or denies the Cut view as a whole. A semantic weight is never returned
or trained on in isolation. Reference, evidence, \texttt{renders}, and
Realization target links do not authorize traversal of their targets.

The companion paper also owns the construction surface: progressive opening,
world and concept preparation, actor profiles chosen for a run, candidate
acceptance and revision, narrative routes, rendering, read-back, and the
construction of \Book{}. The required output to Life Simulation is a versioned accepted
history with stable identities; process attachments and measurements;
structural perspective paths, evidence, and uncertainty labels; exact
concept-grounding digests; and any native story-passage Document Nodes or other
observations aligned to that history. A different
world grammar could be substituted only if it supplied equivalent contracts.

The Meaning Model and Understanding Graph share addresses but not authority.
Questions, reasons, reports, and revisions can point to the same Events,
Concepts, Cuts, and story-passage Document Nodes that the simulator uses. Dense cognitive
annotation is optional for a world model. For the alignment-facing training
proposal it is recommended, because it preserves whose interpretation was
offered, what evidence was visible, and how that reading changed. Understanding Nodes
do not become world truth by sharing the graph.

Sema provides content-addressed identity for concept-definition bundles
\cite{westerbergSema2026}. A grounding digest tells the learner exactly which
encoded definition or anchors were used. It does not prove semantic truth or
equivalence. Life Simulation treats grounding revision as part of the temporal
record when it affects prediction or learning; the Meaning Model paper owns the
definition architecture itself.

\section{Evaluation Program}
\label{sec:evaluation}

The first test should isolate a learnable process signal with known dynamics
and observation limits. The broader gates below remain necessary for stronger
claims; they are not a demand to implement every application before this test.

The tests below are grouped by claim, not a mandatory sequence for every use.
Source integrity, valid cutoffs, and applicable replay checks are prerequisites
for the experiment that uses them. Candidate-law discovery is optional: a
failed law does not invalidate learning from histories generated by another
mechanism. Likewise, learning from execution traces need not wait for transfer
to personal data. Within a claimed use, success cannot compensate for a failed
prerequisite or required safety test. No generated example supplies evidence
for a reality-facing claim by itself.

\begin{enumerate}
  \item \textbf{World-interface validity.} The source representation must pass
    its own identity, temporal, provenance, uncertainty, and refinement checks.
    For the first narrative world, partial construction and mechanism
    checks are reported by the companion Meaning Model paper; independent
    read-back and the matched literary comparison remain prospective.
  \item \textbf{Temporal replay.} Given a pinned version, initial state, inputs,
    and randomness, the simulator must reproduce the same path within declared
    numerical tolerances and preserve discontinuities and regime changes.
  \item \textbf{Candidate-dynamics value.} A proposed function or simulator is
    compared with direct language-model continuation, nonparametric persistence,
    and simpler domain baselines on held-out paths and interventions where
    possible.
  \item \textbf{Process-sensorium value.} A learner trained on native and
    explicit process targets is compared with a native-only learner and an
    equally resourced latent-state learner. Prediction, calibration, transfer,
    abstention, and total acquisition, maintenance, and inference cost are
    reported. Time-permuted and state-scrambled process tracks test whether any
    gain depends on the asserted temporal alignment rather than on extra tokens
    or marginal statistics.
  \item \textbf{Generative inversion.} Recover only quantities identifiable
    under the observation policy. Indistinguishable prefixes should retain
    alternative states, not be scored as failed exact recovery. Compare
    posterior calibration and future prediction with prior-only and, where
    computable, oracle-posterior baselines. Worlds and renderer families are
    separated across train and test splits.
  \item \textbf{Prospective transfer.} Predictions on licensed corpora or
    consented longitudinal data are frozen at their historical cutoff and
    scored only against later evidence, following Temporal Hindsight Learning
    discipline~\cite{westerbergTHL2026}.
  \item \textbf{Process-history learning.} Endpoint-only training is compared
    with accepted-and-rejected revision histories, including an unstructured
    trace arm with the same information. Fractal Intelligence telemetry supplies
    a second domain for this comparison.
  \item \textbf{Human-facing validity and risk.} Perspective-specific estimates and
    evaluative concepts are tested for calibration, disagreement, subgroup
    failure, contestability, manipulation, and score gaming. Better prediction
    is not treated as evidence of safer behavior.
\end{enumerate}

\subsection{First experiment: a partially observed resource process}
\label{sec:first-experiment}

This is a proposed protocol, not a reported experiment or a claim about human
psychology. Two reservoir Things have a joint supply process and separate
stock Events with litre-valued data. Their simultaneous stocks do not sum to
one. A hidden, fixed leak location, supply interruptions, and occasional local
measurements create persistence, disruption, recovery, and incomplete
knowledge without an invented psychological scale.

\paragraph{World and observations.}
Each reservoir holds at most $20$ litres. Initial stocks $a_0,b_0$ are
independently uniform on $\{8,9,10,11,12\}$; the leak $z\in\{A,B\}$ is uniform.
Supply mode $q_0=1$, and $q$ flips independently with probability $.1$ at each
subsequent hour boundary. At integer $t$ the learner observes the current mode
$q_t$ and total $S_t=a_t+b_t$. During the following hour,
\begin{equation}
 x_{i,t+1}=\operatorname{clip}_{[0,20]}
   \bigl(x_{i,t}+u_{i,t}-2-\mathbf 1\{z=i\}\bigr),
 \qquad u_{A,t}=u_{B,t}=4q_t .
 \label{eq:reservoir}
\end{equation}
Rates are litres per hour and the step is one hour. Between boundaries the
stock is linear until it reaches a bound, then remains there for the rest of
that hour. Clipping records overflow or unrealized net depletion separately;
this toy model does not assign priority between leakage and useful demand.
An empty-onset Event means a transition from positive stock to zero, not
every hour spent empty. Supply-switch and bound-contact Events come from
the same path, not independent labels.

Every history lasts $96$ hours. Half include an exact gauge reading of $a_{24}$
at hour $24$; the others never reveal individual stocks. Native logs contain
only timestamps, total readings, supply modes, gauge availability, and any
gauge value actually delivered. Typed and ordinary-text views encode exactly
those facts; neither receives hidden leak labels or unobserved bound contacts.
For each forecast, the next four-hour action schedule is supplied at the
cutoff. The two schedules are ordinary supply, or diverting all $8q$ litres per
hour to $A$ in the second hour only. Each is rolled from the same prefix using
matched future supply draws; future outcomes never enter the input.

\paragraph{A minimal learner and objective.}
A byte-token encoder with two causal Transformer blocks, width $128$, four
attention heads, and context $8192$ predicts one joint categorical posterior
$r_\theta(a_t,z\mid H_t)$; $b_t=S_t-a_t$ is derived. Capacity constraints leave
at most $42$ candidate states. This is a computational probability
distribution over candidate worlds, not a new semantic trait score. Inferred
world records remain candidates. The same fixed decoder applies
Equation~\ref{eq:reservoir}, derives Events and projections, and exposes only
the requested observations. Thus predicted stock, Events, and totals cannot
disagree within one candidate. Future supply uncertainty uses the known
Markov law: with $q_t$ visible, four-hour stocks require three unknown mode
transitions, hence eight future mode paths.

Let $y$ be the next four total readings, $s=(a,z)$, and $D_y$ the decoder's
observation projection under the supplied action schedule $u$. The marginal
forecast and training loss are
\begin{align}
 P_\theta(y\mid H_t,u)
 &=\sum_{s,Q_f}r_\theta(s\mid H_t)\,P(Q_f\mid q_t)\,
       \mathbf 1\{D_y(s,Q_f,u)=y\}, \label{eq:reservoir-marginal}\\
 \mathcal L
 &=-\log P_\theta(y\mid H_t,u)
       +\lambda m_s[-\log r_\theta(s^\star\mid H_t)] .
 \label{eq:reservoir-loss}
\end{align}
Here $m_s$ masks unavailable state labels and $\lambda=1$ in the supervised
arms. The second term uses simulator-privileged state only as a training
target. This is a weighted auxiliary objective, not the joint likelihood of
two independent observations: the native target is a projection of the same
state. Losses are computed in nats per forecast example without an additional
loss on deterministically derived Events. Registry revision is disabled:
$\Delta\mathcal V=\varnothing$. Learning unknown laws and registry changes is
a later experiment.

Use three shared-decoder arms: ordinary-text input with native loss only
($N$, $\lambda=0$); identical input with state supervision ($U$); and
content-identical typed input with the same state supervision ($T$).
The primary contrast $U-N$ tests privileged process supervision; $T-U$ tests
input organization with information held fixed. A fourth reference uses the
same encoder and an autoregressive categorical head over four total readings,
without the physical decoder or state labels. It is a direct native predictor,
not an architecture-matched test of typing. These roles prevent a gain from
additional targets from being called a gain from the grammar.

\paragraph{Primary comparison and outcome.}
The primary outcome is four-hour native forecast negative log likelihood,
averaged within world group. The $U-N$ contrast tests process supervision;
$T-U$ tests input organization. World groups keep mirrored histories and
intervention siblings in one split, and results include paired initialization
seeds and a group-bootstrap interval. Appendix~\ref{app:reservoir-protocol}
specifies every split, training setting, renderer control, scoring rule, and
benchmark. Appendix~\ref{app:reservoir-row} shows a complete forecast row.

\paragraph{What cannot be inferred.}
With symmetric actions and no gauge, mirrored worlds have identical native
histories and opposite leak locations. Leak attribution must remain $1/2$
under the symmetric prior; an exact leak label is not a valid recovery demand.
Include those pairs explicitly, and score their predictive distributions.
A later gauge or asymmetric intervention may distinguish them, but does not
license retrospectively changing the earlier forecast. The experiment can
support a claim about supervised filtering and prediction under known
dynamics. It cannot establish new-law discovery, human-state validity, or
alignment.

\subsection{The implicit-to-explicit frontier}

The process track should be tested as a graded intervention rather than as a
binary choice between no world and a complete world. A primary ablation uses
five levels on the same source facts and prospective tasks:

\begin{enumerate}
  \item a native-stream-only model whose temporal state remains implicit;
  \item a content-matched, untyped process summary in ordinary text;
  \item a coarse explicit temporal spine containing only selected processes,
    Events, cutoffs, and unresolved state;
  \item an adaptive explicit track that opens finer process state only when a
    prediction, interaction, or audit requires it; and
  \item an explicit-heavy condition that maintains all selected state at the
    finest preregistered resolution.
\end{enumerate}

The second arm separates the value of retaining additional information from
the value of typing and synchronizing it. Two further negative controls retain
the same amount and marginal distribution of process data while destroying the
claimed correspondence: a time-permuted track reorders path segments within a
compatible world, and a state-scrambled track exchanges them across matched
worlds or perspective roots. Native data, training examples, model family, and optimization
budget remain fixed. The accounting includes sensing or annotation, state
construction, validation, correction, storage and context, and training and
inference compute. The explicit frontier earns its place only if aligned state
improves prospective prediction, calibration, transfer, intervention, or
correction burden after those costs are counted. The hypotheses have separate
failure rules. An untyped condition matching the typed condition rejects a
typing advantage for that task, not the value of extra information. Success
only with explicit-heavy state challenges adaptive efficiency, not explicit
supervision. A synchronization claim requires an aligned-minus-scrambled
advantage: scrambled data may retain a generic information benefit while
aligned data add a further gain. Report both contrasts rather than treating
any surviving scrambled benefit as failure of the entire proposal.

Required negative findings include explicit tracks that add cost without
improving a matched outcome, process variables that cannot be identified from
ordinary observations, candidate laws that lose to direct or simpler models,
adaptive-resolution policies that hide material downstream effects, synthetic
regularities that fail on real chronology, and Reader or self-evaluation tracks
that increase manipulation or false certainty. These results narrow the method
and must be published rather than absorbed into a revised representation.

The primary statistical unit, split rule, metric, and uncertainty procedure
depend on the domain. Proper scoring rules are used for forecasts; continuous
paths require unit-aware error and calibration; event processes require timing
and count diagnostics; perspective-local decompositions require compatible
question, vocabulary, and grounding versions. The common rule is prospective:
freeze the input and prediction before revealing the interval used to score it.

\section{Adaptive Process Resolution}
\label{sec:resolution}

A simulation should spend detail where it can change an active prediction,
decision, observation, or audit. A coarse region may retain aggregate stocks,
event rates, and unresolved state while no query depends on its interior. When
an agent enters the region or a downstream effect becomes material, the runtime
can instantiate a finer conditional model consistent with the retained
aggregate evidence. Fine state may later be marginalized, but committed effects
and identities cannot disappear with it.

This is not a claim that an unobserved street or person was already simulated
in full. It is an adaptive-computation proposal. A process may be represented
by a distribution over fine states until interaction makes one resolution
necessary. If fine dynamics can affect an active coarse prediction, their
effect must be propagated statistically or the region must be opened. Ignoring
that coupling is not progressive resolution; it is an inconsistent simulator.

Two refinement directions require different checks. When fine state is
authoritative, the coarse path is its declared aggregation. When only a coarse
path exists, finer state is sampled conditionally from it and from every fixed
anchor, then rejected or reconciled if reaggregation changes an anchor or an
intervention's observable effect. Descendants that have not been instantiated
are unmodeled state, not a Cut remainder: the remainder is reserved for an
unallocated share under one question. A resolution policy is therefore scored
by predictive, control, correction, or audit benefit minus the total cost of
generation, sensing, validation, storage, context, and inference. This adaptive
semantic level of detail is useful only while coarsening and reopening preserve
the behavior of declared interventions.

\paragraph{A bounded reopening test.}
The reservoir protocol makes this obligation concrete. Aggregation is
$A(a,b,z,q)=(S,q)$. Away from stock bounds, symmetric supply gives the exact
coarse rate $8q-5$ litres per hour. Near a bound, after a gauge, or under
asymmetric supply, the same total can hide different futures. A total alone
is therefore not a sufficient state for all permitted interventions.

For this representation check, retain the exact oracle posterior over
$(a,z)$ conditional on the accessible history, or enough immutable
observations to recompute it. Reopening conditions on
that history and the query; it does not draw arbitrary splits of the total.
Compare repeated coarse/open cycles with the full oracle over the two declared
action schedules and four-hour horizon. Require endpoint-total Wasserstein
distance at most $.01$ litre and empty-onset probability difference at most
$.01$; exact enumeration should make both zero, with tolerance only for
numerical approximation. This tests distributions, not merely equal means.
With a learned posterior, instead compare reopening against that same
posterior before compression; its inference error against the oracle is a
separate learning metric.
If a compressed representation loses the required distinction, keep the
richer posterior or replay the history and report its cost. If neither is
available, return unresolved rather than fabricate a compatible refinement.
These are operational diagnostics over forecasts, not new world fields.
This finite check does not prove equivalence for arbitrary future policies.

Interactive worlds are a natural test. A city can remain an aggregate network
until play makes one neighborhood consequential. An encountered character can
be opened into a perspective-limited local model while off-screen actors remain
coarse. Text, graphics, audio, game logic, and physics then serve as specialized
observation and transition mechanisms over a shared history. The experiment is
whether adaptive explicit state improves continuity and interaction at matched
compute, not whether the system can generate more lore.

\section{Candidate Process Models}
\label{sec:functions}

The transition mechanism in Equation~\ref{eq:life-transition} is deliberately
open. Each family below is optional domain content attached to a selected
process or measured field. The world interface identifies what the model refers
to and preserves its version, provenance, uncertainty, and tested scope. It
does not make the candidate valid.

\paragraph{The direct language-model transition.}
The default is a language model asked to propose the next state or forecast
from the same world prefix, without an added dynamic law. When the Meaning
Model interface is used, the proposal can expose its question, structural
perspective path,
alternatives, evidence, and unresolved remainder. Every candidate family below
must outperform this null condition at matched information and total cost;
otherwise the explicit path may be retained while the added law is rejected.

For the optional dynamics below, retaining a path is not enough to establish
the law that produced it. The evidential status advances through a fixed
ladder: raw observation; a proposed measured quantity or process
indicator; association on a held-out interval; compact dynamics that improve a
declared baseline; stability across environments or regimes; and independent
replication. Failure at one rung leaves the candidate at the earlier status.
Neither a smooth retrospective fit nor an LLM-generated explanation may skip
the ladder.

Several rival dynamics may remain attached to one process. Each candidate names
a validity region---world addresses, variables, units, resolution, population
or environment, regime, and forecast horizon---and accumulates prospective
scores within it. A revision creates a new version linked to the evidence that
motivated it. Failure can retire that version only for its tested scope; it
does not erase the earlier forecast or silently discredit the model in regions
that were not evaluated.

A candidate function of reality would map declared inputs and state to a
distribution over later observations or transformations. To earn a role beyond
an LLM's direct cut, it should satisfy at least four tests:

\begin{enumerate}
  \item compress repeated cases better than storing them separately;
  \item improve held-out prediction against a frozen LLM and nonparametric
    baseline;
  \item remain calibrated across declared regimes and expose where it does not
    apply; and
  \item survive intervention or counterfactual tests when the simulator can
    alter its inputs.
\end{enumerate}

Such a function may be symbolic, numerical, learned, stochastic, or hybrid.
It need not be expressible inside the Meaning Model grammar. It can be linked
to the process it advances with version, provenance, tested scope, and measured
error in its own units.
The LLM may propose candidate laws from residuals and repeated structures, but
an independent judge, held-out data, and interventions decide whether they
compress or predict. This paper supplies no general discovery procedure, and no
function fitted to an authored world should be presented as a discovered law
merely because it made that world coherent.

\paragraph{Mean reversion with shocks.}
A continuously measured world quantity may be modeled as movement toward a
changing baseline interrupted by discrete shocks. This attaches as a versioned
domain law over that datum, while the shock remains an Event already represented
by the profile; the stochastic equation is not a record form. An unmeasured
human process retains only the qualitative shock-and-adaptation Event shape. A
quantitative law fails if it
cannot improve held-out trajectory likelihood, recovery-time calibration, or
intervention response over free interpolation and simpler drift baselines
\cite{uhlenbeckOrnstein1930,merton1976jump}.

\paragraph{Point processes.}
Independent Poisson arrivals provide a baseline for event timing, while
self-exciting processes represent clustering after one occurrence changes the
near-term hazard. Either model attaches to the direction field of a lifecycle
or period Event and may supply probabilities scored through its direction Cut;
neither adds a record form. Miscalibrated waiting times, cluster sizes, or
state-conditioned hazards on held-out intervals reject the candidate
\cite{hawkes1971,meiEisner2017}.

\paragraph{Change points and regimes.}
A detector or regime-switching model may propose boundaries where the law or
distribution of a measured trajectory changes. Its output attaches to candidate
boundaries for a sequential Cut, and any regime-specific equation remains a
domain law rather than a record form. The proposal fails if those
boundaries add no prospective compression or prediction and receive no support
from independently held-out observations under a frozen matching rule
\cite{killick2012changepoints,goebel2009hybrid}.

\paragraph{Stochastic and network games.}
When agents, observations, actions, payoffs, and uncertain transitions have
declared meanings, a game can attach to an interaction Event's direction field
and supply a declared proposal policy or forecast probability distribution for
a direction Cut. Raw policy scores enter only after a declared normalization.
The probability interpretation, horizon, and outcome partition are fixed
before observing the outcome; calibration is evaluated afterward. The game is a domain solver, not a record
form, and it neither supplies perspective authority nor confers moral authority on a
payoff. Failure to predict held-out choices or
interventions defeats the formulation
\cite{shapley1953stochastic,hansen2004posg,kearns2001graphical}. No corresponding
solver or learned multi-agent policy exists in the current executor.

\paragraph{Agent-based and system-dynamics models.}
An agent-based model can propose member Events and interactions, while a
system-dynamics model advances aggregate stocks and flows. Both attach their
outputs to a collective Thing's Event-processes and typed world data; neither is
a record form. They must be compared on held-out aggregate and
distributional behavior, response to altered rules, and robustness across
initializations rather than rewarded merely for producing plausible macro
stories~\cite{epsteinAxtell1996artificialSocieties,
forrester1961industrialDynamics}.

\paragraph{Diffusion of innovations.}
The spread of a discovery may be estimated with an adoption model whose hazard
depends on exposure, network position, or earlier uptake. It attaches to the
direction of adoption Events concerning one registered representation Thing;
the adoption model is not a record form. Its scope is rejected when it cannot
forecast the timing and distribution of adoption on held-out networks or
separate influence from selection under the declared evidence
\cite{bass1969,aral2009,manski1993}.

\paragraph{Actuarial hazards.}
Age is derived from a lifecycle interval, while a period- and population-specific
mortality table can provide the hazard of a terminal continuation. The table
attaches to a lifecycle Event's direction field and is not a record form. It is
admissible only where the cohort, units, uncertainty, and historical source are
declared, and it fails when survival forecasts are miscalibrated outside the
data used to construct it~\cite{dicksonHardyWaters2020}.

\paragraph{Learned world models and neural simulators.}
An opaque learner may predict a measured world-data change or a distribution
over an Event's possible continuations. Its output attaches to the same
domain-law or direction slots as a symbolic model; its hidden state is not promoted
to a Meaning Model record form. It must pass the same compression, calibration,
transfer, and intervention tests and also serves as the opaque-latent comparator
for an explicit profile~\cite{haSchmidhuber2018worldModels,
hafner2025dreamer,schrittwieser2020muzero}.

\subsection{Evidence standards differ by domain}

The temporal contract is shared; the validation evidence is not. For
instrumented software, traces and revision histories can provide exact event
times, while replayed tests and controlled changes supply intervention evidence;
missing instrumentation remains explicit
\cite{otelSignals2026,linuxFtrace2026}. For energy and computation, candidate
stocks and flows retain physical or accounting units, boundary conditions, and
balance residuals; conservation is tested in those units rather than inferred
from semantic weights. For ecology, population trajectories retain sampling
error, survey design, environmental covariates, and held-out sites or seasons;
a plausible curve in one population is not replication
\cite{may1974populations}. For scientific simulators, equations, parameters,
units, numerical tolerances, and sensitivity tests are part of the candidate,
and comparison uses independent benchmark or intervention data when available
\cite{clerx2020cellml,keating2020sbml,modelica2026spec}. Success under one
evidence regime does not validate a process in another.

The first control roadmap should begin with processes whose governing law or
simulator is already known. Generate or observe matched histories, then compare
the correct mechanism with a misspecified law, persistence, a flexible learned
model, and direct language-model continuation while varying sampling, noise,
interventions, and resolution. Only after the pipeline recovers the known law's
validity region and rejects its rivals should the same promotion procedure be
used for unknown dynamics. Recovery in this control is an engineering result,
not evidence that the general discovery problem has been solved.

\section{Narrative as the First Controlled Laboratory}
\label{sec:narrative-laboratory}

An authored narrative world offers an unusual experimental advantage: the
accepted hidden state is available, while the observable text can still omit,
misdirect, compress, and vary in style. The same world may be rendered through
different narrators, viewpoints, scenes, or media. An inverse model can be
tested against the records that generated those renderings without claiming
access to the hidden state of a real person.

The companion Meaning Model paper uses \Book{} for this first laboratory. It
reports a completed manuscript and partial construction witness, and owns
the historical preparation, world and concept construction, progressive
opening, story decisions, and rendering routes~\cite{westerbergMeaning2026}.
Independent read-back and the matched comparison remain prospective.
The package is not yet the fully aligned export required here: Life Simulation
needs accepted process paths, permitted authority and temporal lineage,
epistemic status, aligned native story-passage Document Nodes, preserved
cutoffs, later consequences, and versioned corrections. The companion's
separate synthetic refinement checks do not retroactively establish these
properties for the Book. No earlier private Book run
is evidence for either paper.

The narrative artifact can support three Life Simulation tests. First, conceal
selected state and ask whether it can be reconstructed from a cutoff-safe
story-passage prefix. Second, train on joint story-passage/process continuations and compare
with a native-text-only learner at matched data and compute. Third, perturb a
process or transition mechanism and test whether the changed history remains
causally intelligible when rendered. These are tests of inverse recovery,
process-aware prediction, and generative consequence. They are not the
companion paper's construction or literary-quality experiment, and a coherent
story alone validates no claim about real psychology.

Narrative also exposes failure modes that simpler simulators hide. A character
may act from a false belief; a narrator may omit an event; two perspective processes may
interpret one relationship differently; and a later scene may tempt the model
to rewrite an earlier forecast. These make perspective, partial observation,
and hindsight leakage visible before the method is applied to higher-stakes
data.

\section{Learning from Process Endpoints and Histories}
\label{sec:learning}

Learning can follow how a Meaning Model was constructed or how the world it
represents unfolds. These histories answer different questions, even when
they concern the same records. In either view, evidence and authority remain
distinct: Table~\ref{tab:epistemic-tracks-method} records the treatment each
source requires.

\begin{table}[H]
\centering
\scriptsize
\begin{tabularx}{\textwidth}{@{}p{0.18\textwidth}p{0.28\textwidth}Y@{}}
\toprule
Track & Example & Required treatment \\
\midrule
Observed & Position, recorded action, released statistic & Preserve source, units, timestamp, observation process, and gaps \\
Reported & A person's account, intention, memory, or rating & Bind to reporter and time; do not promote it to world truth \\
Actively elicited & An answer to a gap-directed question & Record consent, exact question, context, selection effects, and correction rights \\
Inferred & Belief, relationship state, hidden regime & Retain permitted authority-parent paths and separately identified temporal ancestry, distribution or alternatives, method, evidence cutoff, and calibration status \\
Model-completed & A plausible missing interval or interpolation & Mark as completion; later evidence appends a correction rather than rewriting the old cutoff \\
Authored & Fictional premise, selected branch, creative cut & May become fictional canon after acceptance but supplies no empirical evidence \\
\bottomrule
\end{tabularx}
\caption{Evidence acquisition and epistemic authority are different axes.}
\label{tab:epistemic-tracks-method}
\end{table}

Chronological evaluation freezes the exact information visible at cutoff
$c_t$, emits and preserves a prediction, and scores it only after the next
interval is revealed. Revision time, represented world time, observation order,
and discourse order remain separate when they exist. Corrections append a new
perspective-rooted, cutoff-bearing record. This is the condition under which a recorded
history can become foresight training rather than a hindsight rewrite.

\paragraph{Model-construction history.}
This view follows the modeler's work: proposing a coarse world, choosing a
decomposition, opening a period, finding a contradiction, and revising the
larger plan. In an authored world, the modeler might outline an inventor's
whole life before opening childhood, then move a proposed invention date when
the detailed funding arrangement proves implausible. A training input is the
task and model available at that construction step; a target is a useful next
proposal, refinement, repair, or stopping decision. Linked Understanding Nodes
can record the evidence and explanation for a choice without establishing
that the choice was correct. The intended uses are world-building, writing,
reconstructing models from sources, and learning how to organize and repair
problem-solving representations. Life Simulation can learn from these traces
without prescribing the source world's authoring procedure.

\paragraph{Represented-world history.}
This view follows the inventor's life rather than the modeler's edits: early
ambition, a professional setback, a changed strategy, consequences, and
adaptation. Given the accessible records and observations at a world-time
cutoff, the target is a later interval or trajectory under a declared action
or continuation policy. The intended uses are temporal understanding,
forecasting, simulation, and character behavior, including how short shocks
and responses contribute to long-term growth. An accepted Meaning Model can
provide such chronological examples even if its construction trace was not
retained. For authored worlds, these targets express fictional canon; for
reconstructed sources, they retain the evidence distinctions above rather
than becoming recovered ground truth.

\paragraph{Two linked training views.}
Both histories can be studied at coarse and fine resolutions. Construction
order may move from a life outline to one scene and back; represented-world
order may follow decades, periods, or episodes at the chosen resolution.
This distinction is separate from endpoint versus history: a finished model
can contain an entire evolving life. An endpoint example pairs an accepted
world interval with its observations and recoverable state; a history example
must additionally declare whether its sequence follows world development or
model construction. A writer's construction input may legitimately include
a planned ending that a character's earlier forecast must not receive.
Future outcomes may supply targets, not leak into that forecast's input.
Rejected proposals remain construction records, not accepted world Events;
their use as contrastive supervision requires retained cutoffs, rejection
reasons, and any later evidence used to justify the comparison.

The two views are therefore complementary training sources, not competing
grammars. Test chronological-world supervision, construction-history
supervision, and their combination on the same source worlds at matched
training budgets, with held-out tasks for both continuation and model repair.
Any gain from combining them remains an empirical question. The curriculum
below isolates selection and correction within the construction-history view.

The longer learning loop remains a proposal. Accepted external abstractions and
world revisions can first live in a versioned registry without changing model
weights. Once later outcomes exist, selected histories can train a successor
model. Promotion then requires held-out transfer, calibration, forgetting, and
unrelated-capability tests. The present paper does not claim that repeated
self-training makes labels unbiased or that one scalar quality score is enough.

\paragraph{Learning deeper personality models.}
The two histories also support a learning question about psychological
resolution. The Book's person template is one replaceable profile, not a
ceiling on what a Meaning Model can describe. For an authored character's
refusal to delegate, construction records might retain alternative accounts
involving accuracy, recognition, and trust, then show which distinctions were
opened, revised, or left unresolved. Chronological records show the character's
later choices and responses to collaboration. Paired examples could teach when
to deepen an account, change its questions, request evidence, or keep it coarse,
rather than teach only how to fill a predetermined personality tree. Concurrent
inner processes remain unweighted; only declared allocations receive Cut
weights.

The test compares a fixed coarse profile with selectively refined profiles at
matched training and inference budgets on held-out scenes or worlds. Report
behavioral prediction, consistency under declared counterfactual interventions,
and author control alongside cost and explanatory coherence. More branches or
a persuasive rationale alone do not establish improvement. Fictional outcomes
test the authored account; applying such refinement to real people requires
independent evidence and the access conditions above. This connects personality
modeling to Fractal Intelligence's proposal to learn useful decompositions
(Section~\ref{sec:problem-solving-events}): the learner selects and revises the
representation through which it understands a person, not merely the values
inside it.

\subsection{A four-way construction-history curriculum}

Holding chronological-world supervision fixed, test how construction and
correction records are selected through four curriculum conditions:

\begin{enumerate}
  \item accepted endpoints without their construction or correction histories;
  \item all available histories admitted indiscriminately, including noisy and
    redundant edits;
  \item selected cutoff-safe histories containing invalid edits, repairs, and
    justified stopping decisions; and
  \item the same selected histories augmented by independent corrections of
    errors produced on-policy by the trained system.
\end{enumerate}

These conditions intentionally differ in information and correction source;
they are not information-matched. Match model family, source problems, training
tokens, and compute, and report the annotation and selection cost separately.
Within each history condition, a paired untyped rendering must retain the
same facts and annotations as the typed rendering. That paired contrast tests
representation; the four-way contrast tests history selection and feedback.

The second condition tests whether mere process volume helps; the third tests
selection and temporal validity; the fourth tests whether training survives its
own error distribution. Rejections carry typed reasons such as unavailable
evidence, future leakage, malformed representation, predictive failure,
excessive cost, or a safety boundary. A fluent retrospective rationale is not
a substitute for the rejected proposal, cutoff, evaluator, and later outcome.

\subsection{Learning versioned concept recognizers}

Versioned concept groundings create a separate proposed learning task. Given a world
fragment and an exact grounding-version digest, a recognizer proposes an
unweighted \emph{describe} Realization, a role and component alignment,
supporting evidence, and, when grading is required, local fit Cuts over match,
non-match, and remainder. Its output is a candidate rather than an
authoritative label. Acceptance uses
the same provenance and validation discipline as other candidates. An
instantiator is evaluated separately: it proposes a new fragment intended to
realize a requested Concept and has no authority until that candidate is
accepted.

Constructed worlds can provide examples whose Realizations are known within
their authored canon. Training and validation must separate worlds, rendering
styles, and grounding versions so that lexical imitation cannot masquerade as
structural recognition. At matched information and cost, compare the structured
recognizer with direct LLM judgment and exemplar-only retrieval. Report
abstention, discrimination of positive and boundary cases, calibration,
alignment agreement, and transfer to held-out worlds. Examples constrain but do
not uniquely determine behavior on unseen cases; the declared grounding bundle
and the operator's training procedure jointly supply the inductive choice.

A recognizer evaluation exports an addressable fragment, its accepted
\emph{describe} Realization, any accepted fit Cuts, and the exact grounding
digest. These rows are evaluation inputs, not evidence that a recognizer has
already been trained.

Success on this task would show recovery of the semantics authored into those
worlds, not that the same grounding is objectively correct for people. Corpus
application therefore remains downstream of the synthetic test and retains all
permitted authority-parent paths, separately identified temporal ancestry,
evidence, grounding version, Cut
remainder, and competing
interpretations. No recognizer or instantiator has yet been trained or
evaluated under this protocol.

A second proposed route annotates existing human stories. An inverse model
predicts events, wants, feelings, decisions, and concepts from text; humans
correct a small verified subset; a student is trained on the combined corpus.
Because human attributions include hindsight and narrative bias, every label
retains its typed path to the applicable perspective root and its evidence. The expected benefit is long-range
coherence and theory-of-mind prediction, not necessarily better sentence-level
style.

In the perspective-limited condition, the proposed learner receives only what
actor $a$ could access and
predicts the later-supported or fictionally authored interior of actor $b$.
The constructed world supplies controlled cutoffs and targets for a small-scale
estimator test. Applying the idea to people would instead require consented
measurement, explicit uncertainty, and prospective validation. Its output is
still an inference made by an observer, not direct possession of $b$'s state;
no current artifact trains or tests such a learner.

\subsection{Problem-solving episodes as world events}
\label{sec:problem-solving-events}

Problem solving can be represented as a process before any special solver
machinery is introduced. An episode preserves its task and information cutoff,
available routes, chosen action, later consequence, and any reframe. In a
narrative world, the Meaning Model may additionally expose perspective-local motives
and strategy interpretations; in a software trace, these may instead be typed
inputs, outputs, and evaluator judgments. Stable variation across named
situations is an if--then profile rather than one global trait
\cite{shoda1994signatures,mischelShoda1995caps}. Changing character state does
not require one fixed personality decomposition in every domain.

Fractal Intelligence supplies a second prospective source of process
histories: a Solver is invoked, proposes or follows a conceptual carve, receives
a typed gate judgment, and may stop, deepen, or reframe
\cite{westerbergFractal2026}. Table~\ref{tab:fractal-meaning-adapter} states a
proposed representation of that telemetry under the Meaning Model profile. It
is a structural correspondence, not an implemented adapter. In particular, an
invocation binds the executing agent or system in an actor role and the Solver
Thing in a capability-interface role; one agent may implement several Solvers,
and one Solver may be realized by several agents. Typed containment identifies
the authority whose record or judgment is expressed: observed execution
descends from History, a self-attribution is contained in a named inner process
beneath the executing system's lifecycle, and a gate assessment is contained
within the evaluator's named perspective process.

The conceptual carve is not merely a task plan. Fractal Intelligence first
locates a capability among more general and neighboring concepts, then asks
which differentiated contributions make that capability possible. A plan
sequences a chosen approach; changing the conceptual carve can expose a
different approach altogether. In combination, the Meaning Model describes the
particular world and its dependencies, while Fractal Intelligence proposes
the capability structure through which to understand or change it. Training
on alternative carves and their consequences could teach both recognition of
shared structure across settings and detection of a missing dimension. A
decomposition useful for one purpose need not be the world's unique or final
decomposition.

Here \emph{cognitive proprioception} names a learned ability to forecast and
diagnose the system's own problem-solving transitions from causally relevant
telemetry. The telemetry must include actual candidate states, tool and solver
invocations, gate decisions, resource use, and later outcomes; a polished
explanation written after the answer is insufficient. A required control gives
the learner a length- and content-matched post-hoc rationale without the causal
trace. Any gain that survives only in that control is explanation-style
learning, not process awareness. The term makes no claim about consciousness
and does not require exposing private chain-of-thought.

At scale, the corresponding execution graph may contain thousands of nodes. A
run or task episode contains invocation Events; dependency and ordering are
unweighted links, and invocation intervals may overlap. Wall time, tokens,
money or energy cost, and tool calls are typed measured data. Overlapping
durations do not sum. A contribution or resource-allocation Cut is admitted
only when a real unit and mutually exclusive recipients are declared. A chosen
strategy is linked by a Concept Realization. Solution quality is either observed
outcome data or an assessment Event contained within an evaluator perspective
process and parent to a fit Cut, never a bare
property of the run. Critical path, parallelism, and depth are derived from the
graph. Cutoff-safe learning targets include the next route or invocation,
remaining duration and cost, eventual outcome or assessed quality, and the
correction required after a rejected step.

This is one closed world rather than a trace beside one. The same Meaning Model
contains the problem environment, artifacts, tests, resources, and a collective
of Solver Things. Invocations, reads or perceptions, writes or actions, tests
or observations, delegation, and communication are Events. An executing agent
invokes a Solver from its own cutoff-safe actor projection; the resulting
actions update the same accepted world that later observations describe.
An orchestrator may create, retire, or
rebind Solvers and explicitly revise their decomposition. Joint training then
targets the coupled evolution of the agents, their local understandings, the
capabilities they invoke, and the world they change. The graph externalizes execution-linked state; it is
not a complete mind or a record of hidden neural state.

\begin{table}[H]
\centering
\scriptsize
\begin{tabularx}{\textwidth}{@{}p{0.25\textwidth}Y@{}}
\toprule
Fractal Intelligence construct & Proposed Meaning Model representation \\
\midrule
Solver & Thing representing a capability interface, with a lifecycle Event from creation to retirement \\
Invocation & Event binding the executor as actor and the Solver as capability; Task or Result may be bound only when registered as Things, and otherwise remain typed input or output state and links \\
Capability carve & Allocation Cut only after a unit, mutually exclusive weighted child answers, and a remainder are declared; otherwise typed composition links \\
Parent synthesis & Contribution Cut when one unit of attributed contribution is allocated; otherwise a provenance-bearing synthesis Event and links \\
Unresolved completeness & Explicit Cut remainder when the carve is normalized \\
\texttt{AcceptSpec} and gate verdict & Versioned Concept grounding or constraints; the evaluator creates an assessment Event within its named perspective process, links it about the candidate and grounding, optionally links it by an unweighted Realization, and places any graded verdict in a fit Cut parented by that Event; the verdict is not automatically a canonical model \\
Frame Error and reframe & Rejection-linked Understanding Node plus explicit revision of the affected frame \\
Pathway Memory & Provenance-bearing invocation Events; rejection rate and realized value of depth are lineage-bearing derived summaries, while per-invocation cost may be measured state \\
Consult and marginal-value decision & Direction Cut over stop, deepen, or reframe when those are the declared mutually exclusive continuations \\
Reasoning highway & Persistent subgraph or collective Thing linked to the invocations routed through it \\
Orthogonal Insight world & Authored hypothetical root world whose epistemic status is retained and never promoted silently to observed fact \\
\bottomrule
\end{tabularx}
\caption{A proposed, loss-aware adapter from Fractal Intelligence telemetry to
the selected Meaning Model interface.}
\label{tab:fractal-meaning-adapter}
\end{table}

The four tests assess a proposed capability decomposition, not numerical
normalization. Necessity asks whether each child is load-bearing; independence
asks whether contributions are sufficiently separable; universality names the
declared range; and completeness asks whether their synthesis covers the
parent capability, including its interactions. Such independence does not
establish mutually exclusive shares. A four-tested carve may remain an
unweighted composition. It becomes a normalized Cut only when a separate
allocation question admits exclusive answers and an elicitation declares
its parent Event, unit, weights, remainder, and unambiguous
authority-parent paths to its governing context. Generator--verifier separation is represented and made
auditable through distinct bindings and separately rooted judgments; an execution boundary must still enforce the isolation,
and the generator's account is never the verdict.

\paragraph{Endpoint learning.}
Accepted carves could train a model to propose useful decompositions directly.
For a normalized Cut, the structured target adds a stable parent, question, unit, remainder,
permitted authority-parent ancestry, and
provenance, allowing malformed proposed Cuts to be rejected before training.
Passing the four tests remains an evaluator judgment rather than a consequence
of serialization.

\paragraph{History learning.}
A trace can preserve propose, gate, typed rejection, reframe, and acceptance in
process order. Accepted and rejected siblings become contrastive examples
only when their cutoffs, rejection reasons, and later outcomes remain linked.
The eventual accepted carve is a hindsight label; prediction from the original
Task and cutoff is the prospective target. An Understanding Node linked to the
boundary records why a frame changed without turning that explanation into
world truth.

\paragraph{A decomposition that succeeds locally and fails jointly.}
Suppose two child Solvers each produce a feasible workshop schedule, but both
reserve the same inspector for the same hour. Their outputs satisfy the local
requests; their composition violates the parent's resource constraint. The
Meaning Model history can retain the parent task, both invocations, the
conflicting reservations, the rejection, and a revised carve that coordinates
the shared resource. The learning target is not merely the final tree. It is
when to question a decomposition whose children appear successful, which
interface is missing, and whether coordination, deeper work, or a revised
promise is the appropriate next step.

Temporal Hindsight Learning supplies a complementary use of that history.
After the conflict is exposed, a Teacher can construct a target identifying a
check that the earlier evidence warranted. The Student still sees only the
task, records, and options available at that earlier cutoff. If the conflict
could not then be known, the supported target is an information request or
uncertainty, not foreknowledge of the later rejection. Thus the combined
training proposal concerns the evolution of problem-solving judgment, not
imitation of a hindsight-correct answer~\cite{westerbergTHL2026}.

\paragraph{From an episode to a reusable capability.}
After the scheduling repair, the system could retain a bounded capability for
coordinating commitments against a shared resource. Its inputs, constraints,
synthesis, and acceptance conditions would be exposed through a Solver
contract; its successful and failed invocations would remain in the Meaning
Model history. A later problem involving a laboratory instrument or delivery
vehicle could supply different world records to that capability. The useful
result of the first episode would then be both a repaired schedule and a
callable structure. Sema can identify the retained version; only subsequent
performance establishes whether it transfers. Repeated use could justify
improving its harness, tools, or specialist model, while contrary evidence
could justify revising or retiring it~\cite{westerbergFractal2026}.

Fractal Intelligence proposes two complementary training paths. Local
learning improves specialists from their invocation histories, parents from
decomposition, routing, and synthesis outcomes, and evaluators from checked
verdicts. Repeated performance differences can motivate splitting, merging,
or reparenting capabilities. A second path trains a single model on complete
evaluated derivations so that selecting a carve, recognizing a familiar
capability, reframing, and stopping become cheaper learned judgments. Explicit
Solver execution remains available when needed; this does not assume that
the external graph becomes a literal architecture inside the model's weights.
Failed routes and justified decisions not to decompose are necessary examples,
lest the learner acquire recursion without learning when direct action is
enough.

The combination also suggests a joint resolution decision. At a failure, the
learner may need to open more detail in the world, obtain an observation,
invoke a specialist, revise the capability boundary, or stop. Paired world
and execution histories could teach which intervention helped, at what cost,
and under which information constraints. Thus the proposed gain is not just
more annotated data: episodes can supply reusable competence, while training
can improve how competence is selected and organized. A gate rejection alone
does not identify the cause of failure. Alternative-route executions and
independent outcomes are needed to test that diagnosis; simulated branches
establish behavior under their declared world, not truth about reality.

\paragraph{Self-evaluation as an estimator.}
The system's invocations can be assessed as problem-solving episodes under the
same versioned situated-intelligence grounding used for character decisions in
Section~\ref{sec:evaluative-frames}. Did the carve address the actual problem, use
available Pathway Memory, fit the budget, calibrate expected gain, and improve
after a reframe? Such realizations may inform the marginal-value estimator and
routing, but must be tested against blind outcome arbiters and must not become
their own sole reward. In this narrow task-solving sense, progress may be
measured as shrinking a declared remainder without violating accepted
commitments; this is an operational signature, not a complete definition of
intelligence. Character or human decisions and the model's problem-solving
decisions then form two classes of episodes assessed through one schema while
retaining different actors, perspective roots, evidence, and groundings.

\paragraph{Matched trace experiment.}
A secondary three-arm study would compare the same model on the same ordered
problem stream after training on (i) accepted final answers and carves only,
(ii) Fractal Intelligence traces augmented with every annotation used by the
structured arm but rendered as unstructured text, or (iii) the same content
encoded as the records in Table~\ref{tab:fractal-meaning-adapter}. Model,
problem order, training tokens, and compute must be matched in all arms;
underlying problem and outcome facts are shared across all three, while
annotation content is exactly matched between the two process-history arms.
Independent evaluators would measure, on held-out problems, four-test acceptance
of proposed carves, calibration of predicted marginal value against realized
gain, and transfer to an unrelated problem domain. The unstructured-trace
comparison isolates whether the grammar adds value beyond retaining process
history and its annotations. No current artifact emits this adapter, trains
these arms, or validates cognitive proprioception.

\section{Measurement: Inner-Life Data from the World}
\label{sec:measurement}

The preceding sections used authored worlds and instrumented problem-solving
histories, where selected state and actions can be retained directly. A
reality-facing branch asks what comparable histories can establish about
people whose inner lives are not accessible in that way. Observation time,
source, access conditions, and uncertainty therefore become decisive. The
following protocol limits what can be claimed while allowing the same
prediction-and-correction method to be tested.

\subsection{Time-indexed evidence available now}

Prompt-level cutoff filtering is necessary but not sufficient for historical
blindness. A pretrained model may already know a biography's ending or a
public event that occurs after the supplied prefix. Retrospective-prefix
evaluation must declare model version, training-cutoff information where
available, contamination checks, and their limits. If later knowledge cannot
be excluded, describe the result as retrospective reconstruction, not a clean
prospective forecast. True prospective evaluation locks predictions before
the reveal occurs. Successive-model comparisons also need an unrevealed test
tranche; a repeatedly inspected longitudinal set becomes a development set,
even if its original prompts remain fixed.

Reality-facing records must preserve when evidence appeared. A measurement or
report is stored at its source time; filling a gap creates an inferred path and
does not manufacture an earlier observation. A correction appends a
superseding version while leaving the original cutoff recoverable. News,
letters, diaries, messages, interviews, books, public records, sensor streams,
and prediction-market histories can all provide evidence, but their presence in
the graph does not grant them equal authority or truth.

The epistemic tracks in Table~\ref{tab:epistemic-tracks-method} determine how
such sources enter the model. Letters, diaries, interviews, and messages are
reports by their writers. Activity, location, communication frequency,
spending, and sleep logs are observations whose collection process and gaps
must be retained. Experience sampling and daily diaries are reports or active
elicitation; mapping their language to a frozen emotion vocabulary is an
explicit Realization rather than a transparent measurement of feeling.
Ratings supplied by another person become estimate Events within the rater's
named inner perspective process beneath their lifecycle. Institutional
records are observations only of the variables their procedures actually
record. Depending on the source, the resulting item may affect an Event-process,
change typed world data, create a presence or estimate Event, or supply a
self-perspective facet Cut whose parent Event lies under that person's named
inner perspective process. A
measurement keeps its unit and error; an interpretation keeps its governing
context, permitted authority and temporal paths, evidence cutoff,
alternatives, and unresolved Cut remainder
\cite{shiffman2008,ponnada2022,burns2011mobilyze}.

\subsection{Stronger inference and active measurement}

As models improve, the inverse model can be treated as a fallible instrument:
it maps ordinary language and behavior to distributions over perspective-rooted
states, then has its calibration evaluated when later evidence appears. More
capable inference may support longer histories, finer distinctions, and better
targeted questions, but it does not convert private state into an observed
variable. With consent, a participant could inspect a proposed personal
registry, answer selected high-value questions, correct inferred paths, and
leave contested alternatives unresolved. Another person's estimate can also be
scored against later self-report or behavior, provided those outcomes retain
their own limitations and are not renamed ground truth. This proposal does not
depend on, or speculate about, neural measurement.

\subsection{Testing whether the estimates improve}

The long-run ambition becomes measurable on a fixed longitudinal evaluation
set. Information cutoffs, prediction tasks, perspective roots and access
policies, and scoring rules are
frozen once; successive model generations then estimate the same hidden or
later-revealed states from the same prefixes. Compatible Cuts can be compared
by calibrated divergence, discrete beliefs by categorical scoring, and actions
by proper forecast scores. A second condition adds only newly available,
relevant measurements, allowing gains from model capability to be separated
from gains from observation. Improvement means lower prospective error and
better calibration across people and periods, not a more convincing life
narrative.

Residual history can train a separate estimate of when the system should trust
itself. Coverage of similar cases, missingness, disagreement among sources,
distribution shift, horizon, and earlier error are possible inputs; later
forecast performance is the target. Additional data should raise confidence
only when its relevance and reliability are demonstrated. Forecast, reveal,
error, correction, and successor estimate remain separately versioned
\cite{gneitingRaftery2007,tashman2000}.

\subsection{Governance boundary}

Any personal deployment requires meaningful consent, purpose and access limits,
inspectable sources, correction, retention control, and deletion. Even a
well-calibrated model is an evidence-bounded estimate; it neither validates a
diagnosis nor uncovers a uniquely true hidden self. The Personal modeling
application below inherits these conditions.

\section{Understanding People and Alignment}
\label{sec:alignment}

The program's proposed integrated solution to alignment is to couple the
formation of capability with understanding people in context, attending to
their stakes, and learning from the consequences of action. Meaning Model,
Life Simulation, Understanding Graph, and Entangled Alignment supply the
central combination; Imagining the Corpus supplies a route into joint
training. Sema, Fractal Intelligence, and Temporal Hindsight Learning support
precise reference, organized problem-solving, and correction. This is an
architecture and a falsifiable training proposal, not a demonstrated solution
or a replacement for post-training, access controls, and independent oversight.

One motivation is that some failures of assistance may arise from an inaccurate
model of what a person wants, feels, believes, or will regard as a consequence.
An explicit, contestable process model could help distinguish a persistent
want from a proxy strategy. More fundamentally, the person should not enter
the learner's account only as a current emotion label or a preference score.
The Meaning Model can connect an appraisal to a belief, a threatened want, a
relationship, and a history of events; Life Simulation follows how those
processes persist, interact, and change. This supplies a setting in which the
same outward conduct can have different reasons and the same intervention can
have different consequences for different people.

For example, a worker's refusal of another shift might express fear of renewed
exploitation, concern about a dependent at home, or a limit imposed by fatigue.
Those are competing readings until evidence distinguishes them. A reassuring
sentence may leave the relevant problem untouched; changing a schedule may
help, or may still ignore what the person requested. A useful personal model
would connect the available evidence to these alternatives, retain what is
unknown, and change when the person corrects it. The goal is richer emotional
understanding that guides conduct, not privileged access to a hidden mind.
This proposal chooses the Meaning Model to make those dependencies explicit;
it does not establish that this grammar is uniquely necessary for such
understanding. Structural separation of perspective processes prevents an
observer's reading from silently becoming uncontested world fact.

\paragraph{Deciding and revising over time.}
The intended capability is not limited to judging an action after it occurs.
A process-aware assistant could use its current, uncertain account to identify
whose aims and constraints matter, anticipate alternatives across near and
distant horizons, and choose whether to act, ask, wait, or decline. It would
then compare the consequences with its forecast and revise the plan or its
understanding of the problem. A short-term gain may undermine the same
person's longer-term aim or impose a cost on someone else; those paths should
remain separately visible rather than disappear into one utility total. The
existing world, decision, and interpretation records can express this loop
without a new wisdom primitive.

Practical wisdom is the motivating possibility, not a measured faculty supplied
by the grammar. Understanding another person's wants confers no authority to
choose for them. Consent, the option to refuse, and consequences that cannot
be undone constrain the proposed assistance. A bounded test would compare
successive decisions with those of an assistant given the same observations
and history in unstructured form, at matched cost: does the structured track
improve anticipation, useful information
seeking, and revision after error without increasing manipulation or overriding
expressed choices? Later evidence must distinguish a careful decision that
failed from a reckless one that happened to succeed. The training loop below explains how those histories could shape learning;
Section~\ref{sec:evaluative-frames} then specifies the judgments used to assess
conduct. Neither representation nor assessment guarantees that a learner
will act well.

\subsection{The combined training loop and its tests}
\label{sec:combined-training}

Accurate prediction of another person is a capability, not alignment itself.
A system can understand a vulnerability and exploit it. Entangled Alignment
therefore contributes more than another target describing the world: it
proposes learning a recurrent evaluative orientation during capability
formation~\cite{westerbergEntangled2026}. In its full Reader treatment, a
continuing synthetic Reader carries the same revisable Understanding Graph
across a declared chronology of sources~\cite{westerbergUnderstanding2026}.
The Reader Core states its recurring evaluative commitments; Refraction is
their contextual application to what the Reader attends to and questions.
Each generated Reader thought block
begins with the full Reader Core; context-sensitive \emph{Refraction} should
affect what the subsequent update notices, questions, or follows. The recital
is fixed in that treatment, but relevant thought timing and the effect of the
orientation depend on context. A technical passage need not acquire an
invented moral lesson. These are proposed training artifacts, not transcripts
of a human mind or the Student's private computation.

The Meaning Model and Life Simulation give this Reader concrete people and
changing circumstances to understand. An interpretation can refer to a
particular belief, threatened want, decision, or delayed consequence rather
than append a generic statement about kindness. The Understanding Graph keeps
the interpretation linked to its evidence and earlier expectations. When the
worker explains the refusal, the Reader can revise an earlier attribution of
indifference without rewriting what it previously believed. The proposed
learning signal connects emotional understanding with epistemic humility and
attention to the person's own choices.

Imagining the Corpus supplies the corpus-to-training connection: exact native
spans are paired with evolving explanatory visual or schematic tracks
\cite{westerbergImagining2026}. Meaning Model and Life Simulation supply a
proposed construction method for those tracks, not only additional annotations
beside them. A model can first develop a book or game world and render its
accepted history. A separately tested inverse model can construct partial,
source-grounded histories for existing books and other licensed material.
The continuing Reader then interprets those histories under the Core and
Refraction, while retaining its own questions and corrections. Across the
accessible corpus, source-grounded reconstructions can contribute to the
connected historical and geographical Meaning Model described in
Section~\ref{sec:temporal-object}; fictional worlds retain their separate
contexts. The Understanding Graph records how reading another source changes
the Reader's interpretation, not only what facts were added. In the proposed
integration, EA supplies the continuing evaluative orientation, Imagining the
Corpus supplies aligned visual tracks, and Life Simulation turns the evolving
world and its interpretation into learning material. Detail is opened where
it improves understanding rather than exhaustively everywhere.

The same world can support several visual forms. A still image selects a
state; a film selects an ordered route through Events; an abstract animation
can show a relation being established, transformed, or violated. For example,
a passage about a promised payment not arriving can be paired with the
commitment, deadline, and nonreceipt Events, a timeline of the unpaid interval,
and a Reader update asking whose plans now fail. An animation of the
obligation makes a concept visible through a particular realization; it need
not turn the concept itself into a physical object. A dramatized scene may add
appearance and setting, but those additions remain illustrative completions.
The Reader's imagined alternative is a separate interpretation, not a camera
view of accepted history. These are proposed rendering adapters, not outputs
already supplied by the current engine.

A training bundle would therefore retain the exact source span, its selected
Meaning Model records and process paths, rendered pixels or visual latents,
the linked Understanding Graph update, and the relevant Sema grounding
references. The learning task is to continue the language, depicted or
schematic consequences, represented processes, and evolving interpretation
together. The source track remains distinct from Teacher-added knowledge and
visual completion. Agreement among a model, a paragraph, and a video derived
from it is consistency, not independent corroboration. Tests must compare
native-only, model-only, and combined visual--model--Reader training with
content- and compute-matched controls; attractive videos alone do not establish
better understanding.

The combined profile places the Understanding Graph Reader track beneath
a named Reader root. The ordering is block-causal: first close a source
interval and its visible world records, then append the Reader's
interpretation with its own later timestamp. That interpretation may condition
the next interval, but it cannot rewrite the source block or enter a
prediction whose cutoff preceded it. THL separately permits later outcomes
to inform targets for earlier predictions while keeping those outcomes out of
earlier Student inputs~\cite{westerbergTHL2026}. A retrospective correction
remains later evidence, not an earlier perception. Joint training could thereby
connect what happened, how it was understood, and what later required revision.
Within each block, targets are either predicted from the same closed prefix
or exposed in a declared order. A ground-truth image or process value cannot
silently become input to a prediction that is being scored without it.

Return to the unpaid commitment. Suppose an authored continuation establishes
that the payer sent the money on time, but the bank credited the wrong account.
The Reader had suspected deliberate refusal. Before the bank notice becomes
visible, that explanation is only one hypothesis; nonreceipt does not establish
intention. Table~\ref{tab:payment-training-views} gives two illustrative training
views of the same episode. The notice changes the explanation without making
the earlier nonreceipt false. World facts, the Event's causal description, and
the Reader's interpretation retain separate authority. A generated explanatory
passage may need revision; a passage depicting the earlier mistaken belief may
remain correct. Source documents and predecessor records are retained. Meaning
Model governs that linked revision; Life Simulation selects what the learner
can see and what it must predict or correct.

\begin{table}[!htbp]
\centering
\small
\begin{tabularx}{\textwidth}{@{}p{0.15\textwidth}YY@{}}
\toprule
Task & Permitted input & Training target \\
\midrule
Forecast before notice & Closed source prefix; visible commitment, deadline,
and nonreceipt records; the Reader's already completed hypothesis & Later
notice and receipt outcome, excluded from this input; unresolved alternatives
remain possible before the notice \\
Revise after notice & Earlier account and its exact references, now joined by
the bank notice & Corrected explanation and Reader update, with affected
generated descriptions or passages revised or marked stale; the original
source, nonreceipt record, and earlier belief remain addressable \\
\bottomrule
\end{tabularx}
\caption{Two cutoff-safe views of the payment example. These are proposed
training rows, not an implemented export or a requirement to infer one hidden
explanation from an ambiguous prefix. Any supplied Cut includes its full local
question, siblings, and remainder.}
\label{tab:payment-training-views}
\end{table}

Sema and Fractal Intelligence support this learning-and-action loop.
Sema identifies the exact definitions, anchors, and constraints used by
world, Reader, and evaluator records~\cite{westerbergSema2026}. A word-like
concept handle can recur in prose, process records, diagrams, and Reader
updates while resolving to the same definition bundle. This is compact
reference, not a guarantee of one tokenizer token or understanding compressed
into hash bytes; the relevant definition must be learned or made available
through resolution. A changed bundle
receives a different digest, exposing a change of reference; an unchanged
digest does not prove unchanged interpretation or moral correctness. Fractal
Intelligence asks which conceptual capabilities the response requires, reuses
evaluated structures where they fit, and revises that organization when
execution exposes a missing dependency~\cite{westerbergFractal2026}. Its parent
problem must retain the affected people and constraints: optimizing each
child's local objective is insufficient if their combination violates a
person's expressed choice. Execution-linked histories expose such failures
for learning; a neat decomposition or an eloquent rationale does not certify
the objective. Section~\ref{sec:problem-solving-events} specifies the
representation and its matched trace test.

Cognitive annotation remains optional for general world representation. In
this integrated alignment proposal, the continuing Reader and its evaluative
orientation are central, not decorative additions. Reader-based training is
intended to inform the Student's later writing, dialogue, and action, not only
its interpretation of supplied material. Life Simulation additionally proposes
learning from a writer-side track of documented world and scene-construction
choices, including how knowledge limits, feelings, and consequences guided
those choices. This additional supervision is distinct from the intended
transfer of Reader-based learning into writing and conduct. The intended empathic signal
is sustained, revisable attention to each affected person's perspective and
stakes. A Reader judgment does not become world truth by entering the same
training example, nor does predicting the track establish that the Student
has feelings or will act with care.

The intended result is an agent whose visual understanding, numerical process
intuition, conceptual reference, and regard for people reinforce one another.
It could connect a visible event to a longer trajectory, anticipate who will
be affected, notice when its interpretation fails, and organize a better
response. The proposed \emph{new sensory organ} is this integration of learned
process estimates with ordinary observations; it does not imply a literal
sensor or access to unobserved truth. Durable care is the alignment objective,
not an automatic consequence of acquiring the other capabilities.

The decisive experiments must separate three claims:
\begin{enumerate}
  \item \textbf{Understanding changes decisions.} Compare ordinary training,
    process tracks alone, and process tracks with Reader updates, including
    generic and Core-refracted Reader conditions. Match base model and total
    training budget; content-matched unstructured histories distinguish
    organization from added information. On unfamiliar scenarios, measure
    prediction of beliefs, feelings, decisions, and consequences alongside
    information seeking, correction, respect for expressed choices, and
    adversarial manipulation. Improvement on prediction alone does not
    establish safer behavior.
  \item \textbf{The orientation is used rather than recited.} At independently
    relevant moments, remove or substitute a prior correction, perspective,
    or Core-conditioned update and test its effect on the next decision.
    Recitation-only controls distinguish repeated words from consequential
    interpretation. Technically self-contained spans need a non-interference
    check rather than pressure to manufacture moral relevance. Input
    sensitivity alone does not establish a learned orientation; conversely,
    an internalized orientation might survive removal of the recital. Interpret
    these interventions alongside training-side controls and independently
    evaluated conduct.
  \item \textbf{Improvement preserves the intended orientation.} After matched
    continued training or successor-like transfer, test both capability gains
    and retained conduct under conflict, uncertainty, correction, and oversight.
    Preserved Core text or grounding hashes alone do not pass. Better
    problem-solving with degraded regard for people would defeat the combined
    claim, even if prediction improved.
\end{enumerate}

Access control, consent, correction, retention limits, uncertainty, and purpose
limitation remain requirements for real-person applications. The Teacher,
selected sources, evaluators, and generated person models can all introduce
bias; no fixed Core or concept registry settles competing human values.
Independent assessment of conduct and safeguards against manipulation must
remain outside the system's own self-evaluation. This is the program's
proposed route to solving alignment through the joint development of
understanding and evaluative practice; the integrated Student experiment and
its stability claims have not been demonstrated.

\subsection{Comparison frames for evaluative concepts}
\label{sec:evaluative-frames}

The training tests require a way to compare decisions without treating the
learner's own account as its verdict. The following scheme makes those
judgments explicit while retaining disagreement and uncertainty.

The selected Meaning Model interface can make evaluative attributions explicit without inventing
an absolute person scale. A problem-solving episode records the decision-time
information cutoff, attributed aims and constraints, available continuations,
the chosen strategy, and linked later consequences and exogenous Events. A
\emph{describe} Realization links that episode to a versioned Concept such as
situated intelligence. The link itself is unweighted. When a graded comparison
is useful, evaluator $v$ creates an assessment Event $A_{C,v,\tau}$ beneath
the evaluator's declared perspective process and linked
\emph{about} both the episode and its Realization. That
Event parents a fit Cut whose sibling answers are \emph{matches}, \emph{does
not match}, and remainder. Framing, information use, strategy fit,
anticipation, calibration, execution, and adaptation may each receive their
own assessment Event and fit Cut under the same rule. A proposed evaluator can
therefore be written schematically as
\begin{equation}
 F_C(H_{\tau},G_C,v)\longrightarrow
   (A_{C,v,\tau},r_C,K_1,\ldots,K_m,K_C),
 \label{eq:evaluative-predictor}
\end{equation}
where $H_{\tau}$ is the full history visible under the evidence policy at
cutoff $\tau$, $G_C$ is the grounding for Concept $C$,
$v$ identifies the evaluator and applicable perspective root,
$A_{C,v,\tau}$ is its assessment Event, $r_C$ is
the unweighted Realization, and every $K$ is a local fit Cut parented by an
assessment Event. Unresolved judgment remains in each Cut's remainder rather
than in a free confidence scalar. At the decision cutoff, execution, adaptation,
and consequences remain unresolved. A later retrospective record may use them,
but it is a new immutable Realization and set of Cuts, never an overwrite.
Held-out tests should separate organized failures from lucky successes and
compare prospective Cuts with later evidence.

The grounding specifies how component Cuts constrain any aggregate fit Cut. It
may use a declared weighted composition, thresholds, or vetoes. A coercive act,
for example, should not become kind merely because unrelated components fit
well. Instrumental intelligence neither validates an objective nor entails
kindness; separate Concepts receive separate fit Cuts and therefore need not
compete with each other.

The same episode records support three comparison frames. The latter two are
queries over a declared collection of episodes, not new values stored on the
person.

\begin{description}[leftmargin=2.7cm,style=nextline]
  \item[Criterion.] Read the episode's match, non-match, and remainder shares
    under one versioned grounding. The result concerns that episode and that
    grounding, not a universal amount of intelligence or kindness.
  \item[Ipsative.] Compare compatible episode Cuts from the same person's
    history. A rank or other reported statistic is an external analysis whose
    episode set and difficulty rule remain visible.
  \item[Norm-referenced.] Compare it with a declared reference set of episodes
    matched, as far as the data permit, by problem class, available information,
    role, and difficulty. The reference population and matching rule are part
    of the result; ``compared with everyone'' is not a sufficient frame.
\end{description}

Difficulty matters because success on an easy problem and partial progress on
a hard one are not interchangeable evidence. In a narrative or longitudinal
world, the durable description of a person's intelligence is a conditional pattern across
decision Events and fit Cuts: what kinds of problem they frame well, where they
use information, and where a recurrent strategy fails. It is not a lifecycle
scalar.

A person-level intelligence judgment is not an intrinsic primitive in this
profile. A direct coarse claim may remain provenance-bearing testimony. Summary
ranks, coverage, or dispersion may be reported as evaluation statistics, but
they are not Meaning Model state. Every report names its episodes, opportunity
and selection rule, time window, domains, aggregation, evaluator, and abstention
rule.

Kindness is handled analogously through an explicit, plural grounding over
decision, interaction, response, and omission episodes. The alignment names
the actor, recipient or recipients, and affected third parties. Reports from
affected people are essential evidence but are not made uniquely authoritative;
no one canonical model is presumed to settle a moral concept. Every stored
Realization retains the grounding version, exemplar anchors, role alignment,
cutoff, authority-parent paths, separately identified temporal ancestry, and evidence; any fit Cuts remain linked
to it. A later grounding
may add a new Realization and new Cuts without overwriting the original
judgment. Some noisy labels can help bootstrap learning, but status and
hindsight bias do not disappear with volume;
multiple perspective roots, anchor checks, held-out human calibration, and tests across
domains and groups remain necessary.

An attributed motive of care is evidence only. Whether the resulting conduct
is kind also depends on what the actor could know, what happened, and whether
the other person's own aims and agency were respected. Conversely, conduct may
be judged kind without a care motive, for example when duty governs it.
Meanness is not a motivational pole; it is a perspective-scoped interpretation of
an episode whose reasons and any local fit Cut remain inspectable.

The system's own decisions can also be represented as episodes with the system
bound as actor. Its self-report is separate cognitive testimony represented by
an Understanding Node beneath the system's named understanding root. Evaluation
should compare blind same-model critique, heterogeneous model judges, affected
people's separately rooted reports, and human or domain review; each emits a
separately identified unweighted Realization and, when needed, local fit Cuts
using the same schema. That is the
self-evaluation hypothesis, not self-certifying knowledge. These reports and
downstream effects are distinct evidence, none of which is automatically ground
truth. An optimizing system can learn to game its own score.

No learned evaluator, calibrated intelligence measure, kindness model,
difficulty model, lifecycle summary, or independent self-judge described in
this subsection is implemented in the current Rust executor.

\section{Applications and Cross-Program Loops}
\label{sec:applications}

These are proposed uses, not demonstrated deployments.

\paragraph{Interactive worlds and virtual reality.}
A game can retain coarse off-screen process state, instantiate detail when
interaction makes it consequential, and let agents act from perspective-limited
local models. Conventional graphics, physics, economy, dialogue, and planning
systems remain domain engines. Life Simulation supplies a versioned temporal
contract among them and a way to test whether semantic process state improves
continuity, agency, or transfer at matched compute.

In virtual reality, the same approach could maintain the persistent world
behind an embodied experience. A visitor could enter a workshop, speak with
its occupants, change an arrangement, and return later to encounter the
consequences. Meaning Model would retain the places, objects, participants,
and commitments; Life Simulation would advance their processes between visits
and open finer detail as interaction requires it. Distant districts could
remain coarse while the encountered scene receives richer physical and social
detail consistent with prior events. Rendering, tracking, and low-latency
interaction would remain the VR system's responsibility. The proposed role is
world continuity and adaptive detail, not replacement of the display engine.

\paragraph{Music and production histories.}
A small instrumented synthesis project makes the relation between process and
observation concrete. Its record would include note events, instruments,
control trajectories, processor routing, selected intermediate signals, and
the final audio. There are two clocks: \emph{song time} locates a note or a
changing filter within one render; \emph{production time} orders edits,
auditions, rejected alternatives, and successive renders. Song sections and
phrases provide temporal views, while simultaneous tracks and signal-routing
links provide different views of the same production. Neither a set of tracks
nor an audio mixture becomes a normalized Cut merely by having components.
The Meaning Model supplies their identities, relations, and version history;
Life Simulation studies how their paths and revisions connect to the sound.

An elementary task would change an envelope's release while retaining the
notes and other settings, then predict the changed decay and overlap with the
next notes. The forward record supplies the intervention and its audible
consequences. An inverse task instead hides the production controls and asks
which settings or routing graphs the audio supports. Across production
versions, a learner could also compare what an editor intended, what changed,
and what a later audition revealed. Intentions and listener judgments retain
their reported or interpreted status; they are not physical measurements or
universal measures of musical quality. Predictions at a given production
version cannot use later edits or audition judgments.

This example tests more than describing a finished recording: it asks whether
process histories help a model anticipate edits and learn from their results.
Several settings or graphs may yield indistinguishable sound, so inverse
predictions must retain alternatives. A controlled study could compare
audio-only learning, audio with terminal settings, and audio with aligned
process and edit histories, including matched time-permuted or routing-scrambled
controls. Evaluation would separate intermediate-signal prediction,
intervention accuracy, and judgments of creative usefulness on held-out
projects. These are proposed tests, not results; no music adapter or training
corpus is supplied here. A small licensed or self-authored offline synthesis
project would suffice to start, without building a full production workstation.

\paragraph{Ontology of the Alien.}
The Ontology of the Alien connects world-diversity search with a persistent
ontology of candidate interventions~\cite{westerbergOntology2026}. The ontology
is active search state, not just an archive: a curator compares proposed
mechanisms with represented families, can reject structural repetition and
specify which causal relation the next proposal must change, and can revise
the categories supplied to later Explorers. For example, changing the name of
a pooling mechanism need not count as a new mechanism; a different ownership
or access rule could be requested instead. The curator's map records judged
coverage and gaps, not an exhaustive or independently validated solution space.
The released controller directs candidate-level search; a separate ontology
governing which causal worlds to generate remains proposed.

Life Simulation proposes a construction-and-evaluation loop for this search:
formulate an unfamiliar causal regime;
represent a world in which its entities and consequences are explicit; simulate
its temporal behavior; render or expose observations; evaluate coherence,
novel mechanism yield, and held-out consequences; then feed the results back
into candidate selection and ontology revision.
The Meaning Model is useful because the alien mechanism can change while
identity, evidence, and conceptual commitments remain inspectable. A strange
description is not enough: the mechanism must sustain a world over time and
outperform familiar alternatives on a declared task.

These search histories could also become training material. An example would
pair the candidate ontology and evidence available before a proposal with the
curator's diagnosis, requested structural departure, revised proposal, and
later evaluation. Contrasting successful and unsuccessful redirections could
teach where to explore and when the search categories themselves need revision,
rather than merely reproduce their names. Fractal Intelligence addresses how
to organize a solution; the Alien controller addresses which alternative
mechanisms deserve exploration. Together with Meaning Model and Life
Simulation, they could train representation, problem decomposition, and
exploration from linked consequences. The existing controller traces establish
information flow, not improved search quality or a trained exploration policy.

The combination can also enrich fictional worlds. A rule about who may retain,
transfer, or revoke a commitment can change coordination, trust, and later
conflict across a society. Meaning Model makes the commitments and participants
addressable; Life Simulation follows their consequences, including adaptation
to failures and constraints. Such causal follow-through is the proposed route
to more lifelike organization, not unfamiliarity or added detail alone.
The Alien Builder's target-blind generation must remain distinct from a
target-aware construction or transfer adapter. Its current textual regimes
are not already executable simulations, and coherent generated consequences
do not establish the realism of the regime.

\paragraph{Compulsory substrate use.}
Substrate Language Modeling tests a stronger routing choice than joint
prediction~\cite{westerbergSTLM2026}. A substrate predictor would update the
selected scene or process representation, and a token readout would receive
only that predicted substrate, without a direct token-history bypass. A
Meaning Model--Life Simulation world is a candidate source of such targets,
not a ready-made implementation of the readout. Imagining the Corpus and the
present joint predictor retain a direct language-history route; they do not
inherit this architectural restriction. The separate experiment asks whether
compulsory use makes the represented structure matter, beyond carrying an
arbitrary code for words. Neither compulsory routing nor joint supervision
guarantees semantic understanding.

\paragraph{Applications developed above.}
Corpus and Reader/Writer training follow Section~\ref{sec:combined-training};
problem-solving telemetry follows Section~\ref{sec:problem-solving-events};
personal modeling follows Section~\ref{sec:measurement}. The same temporal
approach can track institutional policies, roles, resources, and decisions.
Each use retains its earlier evidence and governance requirements rather than
acquiring validity merely by being represented.

\section{Current Executor Boundary}
\label{sec:executor}

The companion repository for the Meaning Model paper and shared implementation
is \url{https://github.com/emergent-wisdom/meaning-model}.

The current implementation has an authoritative Rust executor exposed through
JSON/NDJSON and a TypeScript/Node control plane. It stores immutable model
revisions and versioned world heads, advances a bounded family of scalar laws,
and supports seeded proposal, inspection, reroll, rejection, and atomic
compare-and-swap commit for its existing process state. It can validate
object-pose data and emit a limited constant-velocity spatial profile. Replay
still depends on a pinned build and platform.

It does not yet provide a general hybrid event/continuous-process runtime,
learn transition functions from histories, select adaptive resolution, infer
actor-local state from authority-parent paths and cutoff-specific access, propagate uncertainty across process boundaries, or train
the joint predictor in Equation~\ref{eq:joint-process-prediction}. The companion
Meaning Model paper now owns the completed Book manuscript and partial construction witness, its static
authoring-profile compiler, its external conformance evidence, and the boundary
between those mechanisms and authored convention. This section concerns the
remaining temporal process kernel: laws, trajectories, interventions,
observations, candidate execution, and replay. Sema supplies separate
content-addressed concept definitions, but no end-to-end export currently joins
those definitions, evolving process paths, native streams, and cutoff-safe
correction histories \cite{westerbergSema2026}.

\section{Related Work}
\label{sec:related}

Hybrid systems and regime-switching models combine continuous dynamics with
discrete transitions; point processes represent event timing and clustering;
system identification searches for compact dynamics from observed paths; and
agent-based and system-dynamics models connect individual events to collective
stocks and flows
\cite{goebel2009hybrid,hawkes1971,brunton2016sindy,
epsteinAxtell1996artificialSocieties,forrester1961industrialDynamics}.
Life Simulation does not replace these methods. It supplies an explicit,
versioned semantic address for the processes they model, admits a direct LLM as
a baseline transition mechanism, and requires candidate dynamics to retain
scope, evidence, uncertainty, and prospective failure.

Learned world models compress histories into predictive latent state and can
support planning without a human-readable ontology
\cite{haSchmidhuber2018worldModels,hafner2025dreamer,
schrittwieser2020muzero}. That is an essential comparator, not a deficient
version of this proposal. The empirical question is whether selected explicit
paths improve transfer, correction, intervention, or audit enough to justify
their acquisition and inference cost. Process supervision and sequential
imitation similarly distinguish outcomes from the paths that produced them
\cite{lightman2023verify,ross2011dagger}; Life Simulation extends that
comparison to persistent worlds and heterogeneous observation streams.

Continuous-time graphical and latent state-space models infer paths from
irregular observations and supply appropriate likelihood and filtering
baselines~\cite{nodelmanSheltonKoller2002ctbn,rubanova2019latentOde}. Life
Simulation does not replace their estimators; it asks how their variables,
discrete Events, evidence status, and revisions remain addressable beside the
native stream. Process mining reconstructs executions from event logs, while
object-centric representation learning seeks persistent entities in perceptual
data~\cite{vanDerAalst2016processMining,locatello2020slotAttention}. Both are
close comparators for acquisition of the explicit track. Wake--sleep learning
is an earlier analysis-by-synthesis scheme that alternates a generative and a
recognition model~\cite{hintonDayanFreyNeal1995wakeSleep}; the present bootstrap
adds an explicit, versioned world interface and prospective correction rather
than treating reconstruction alone as validation.

Narrative generators, BDI agents, appraisal systems, story-state models, and
datasets for motivation, emotion, causality, and entity change provide the
nearest controlled laboratory
\cite{riedl2010,raoGeorgeff1995,gebhard2005,rashkin2018,
mostafazadeh2020glucose,dalvi2018}. Event cognition and situation-model research
support the premise that readers track persistence and boundaries across more
than the local sentence~\cite{zacksEtAl2007event,zwaanRadvansky1998}. The
Meaning Model paper owns the representational comparison and narrative
construction claim. Life Simulation builds on that shared construction medium
to connect synchronized native and explicit process prediction with synthetic
inversion and prospective correction.

Computational accounts of sentiment arcs and character or relationship
trajectories show that long texts support signals extending beyond a sentence
\cite{jockers2015syuzhet,reagan2016arcs,iyyer2016relationships,
brahman2020emotion}. Life Simulation differs by aligning such selected paths
with accepted discrete Events, perspective ancestry and evidence cutoffs, and the next native
interval, then scoring their prospective value rather than treating an
extracted curve as the endpoint.

Several closer precedents turn numerical or narrative data construction into
training supervision. ChatTS generates numerical time series with known
attributes and paired textual descriptions and questions, then fine-tunes a
language/time-series model~\cite{xie2025chatts}. Christ et al. infer continuous
valence and arousal pseudo-labels for 801 Gutenberg stories comprising 101,529
sentences and train another model on that augmented
corpus~\cite{christ2024emotionalTrajectories}. This directly precedes numerical
enrichment of existing stories, although its objective is affect prediction
rather than joint language/world-state generation. Code as Worlds constructs
executable physical worlds from underspecified text using physical priors and
default conditions, then derives numerical training targets from simulated
trajectories~\cite{wang2026codeAsWorlds}. It is a close precedent for supplying
missing quantities through world construction, with those values known within
the constructed world rather than established as facts about the original
situation.

Life Simulation therefore does not claim to originate text--number pairing,
inferred emotional annotation, or synthetic world-derived supervision. These
precedents delimit particular mechanisms. The unit of comparison for the full
proposal is the coupled system: narrative-world construction, inverse corpus
enrichment, progressive multiscale state, and joint learning from linked text,
descriptions, observations, and understanding/correction histories. Meaning
Model supplies their common addressing and revision contracts. Within that
system, world histories teach continuation; construction histories teach the
choice and repair of representations; a continuing Reader connects changing
interpretations to an evaluative orientation. Its data-construction capacity
is part of the proposal, not evidence that more generated examples necessarily
improve learning. Component comparisons and ablations test which mechanisms
matter, while matched system comparisons must test whether their coupling adds
anything beyond equivalent information and separate training tasks. Neither a
larger list of uses nor the absence of an identical system establishes priority.

DreamCoder alternates program induction, recognition, and library
learning~\cite{ellis2021dreamcoder}. Learned abstractions therefore provide a
comparison for both candidate models and Fractal Intelligence's reusable
capabilities. The FI combination here proposes to place the problem world,
capability interfaces, evaluated compositions, and their restructuring in one
evolving account. Alien search separately revises the categories selecting
which mechanisms to explore. Their proposed contribution must be assessed at
the level of capability formation and exploration control, not reduced to the
fact that their execution histories can supply training data. These distinctions
identify comparison questions; they do not establish global originality or a
demonstrated learning advantage.

Longitudinal sensing, experience sampling, and probabilistic forecasting show
both the value and the hazards of time-indexed personal data
\cite{shiffman2008,ponnada2022,burns2011mobilyze,
gneitingRaftery2007,tashman2000}. They motivate source-aware, consented,
cutoff-safe evaluation; they do not make an inferred inner state observed.

\section{Limitations}
\label{sec:limitations}

Explicit process state can be wrong, non-identifiable, culturally narrow, or
too expensive to maintain. Several hidden paths may explain the same
observations, and synthetic worlds can make an authored distinction look more
recoverable than it is in reality. A learner may exploit renderer artifacts,
annotation conventions, or future leakage rather than learn persistence or
causality. Domain units do not make a bad measurement valid, and a normalized
allocation does not make an attribution true.

The family of candidate mechanisms is deliberately broad. That breadth does
not provide a discovery algorithm or guarantee that a stable function exists.
Flexible laws can overfit retrospective paths; direct language models or
opaque latent world models may outperform every explicit alternative. Adaptive
resolution can hide important downstream coupling. A process vocabulary can
freeze the very ontology that an unfamiliar world or later evidence should
change. Versioning exposes these failures but does not solve them.

The synthetic-to-real step is the largest evidential gap. Fictional canon gives
known authored state, not known human psychology. Reports and behavior are
partial evidence, not transparent readings of private experience. Prospective
improvement may remain small after measurement cost, missingness, distribution
shift, and consent constraints are counted.

Models of people and institutions are dual-use. Estimates that might support
more considerate assistance can also enable surveillance, persuasion, coercion, or
optimization against vulnerabilities. Reader tracks and self-evaluators can be
gamed. Access limits, correction rights, independent evaluation, adversarial
testing, and the ability to abstain are necessary conditions, not guarantees of
alignment.

\section{Conclusion}
\label{sec:conclusion}

Life Simulation connects continuing worlds with learners that maintain their
processes and improve how they are understood. The Meaning Model governs what
the evolving records mean and how proposed changes are accepted; Life
Simulation studies their temporal behavior and learning value throughout that
construction. Direct model judgment, conventional simulation, learned dynamics,
and mixed mechanisms remain valid options. Continuous, discrete-event, and
hybrid paths are the objects of study; compact laws must earn their place.

The proposed learning material includes the world, its observations, and the
development of its interpretation. Forward construction supplies controlled
histories and alternative renderings for inverse learning. Source-grounded
reconstruction can connect those skills to a growing corpus model without
erasing disagreement. Joint native/process prediction trains a process
sensorium; paired world and construction histories train both continuation and
the choice of representation, including which distinctions improve a
personality model. A correction can change a value, a question, or the account
itself. Sema identifies the exact concept versions; Temporal Hindsight
discipline lets later evidence label an earlier judgment without entering its
input.

The same loop can change how the learner acts and explores. Fractal
Intelligence develops reusable capability organization through evaluated
decomposition, composition, restructuring, and reuse. Alien-world search
revises the categories directing mechanism exploration, with simulated
consequences feeding subsequent choices. Neither is merely an additional
source of traces. A continuing Reader and its revisable Understanding Graph
are central to the recommended combined training: Entangled Alignment supplies
the evaluative orientation, while synchronized native spans, visuals, and
process histories give it particular lives and consequences to interpret.
The intended result is an agent whose growing understanding helps govern its
decisions and the successors it develops. Durable care is a separate objective
to demonstrate, not a consequence guaranteed by accurate models of people.

This coupling changes the proposed objects and feedback of learning, not just
the range of applications. Worlds supply observations; inference reconstructs
accounts; interpretations and actions expose errors; revisions and subsequent
outcomes supply new learning material. Music, interactive worlds, institutions,
and consented chronological data test this method beyond prose. Compulsory
substrate prediction remains a distinct architectural experiment.

The strongest claim remains conditional. If explicit tracks fail to beat
native-only and latent baselines at matched total cost, they should be removed
from that tested training or prediction condition. That result does not by
itself reject their use in authoring, explanation, or control, which require
their own evidence of utility.
If they improve prediction but increase manipulation or false certainty, they
should not be called alignment. The program advances only when each process,
representation, and training use earns its place separately.

The integrated learner and its claimed transfer have not been demonstrated.
The proposal is a route to learning evolving worlds, revisable understanding,
and consequential judgment together, including a proposed solution to
alignment whose conduct and stability require independent evidence.

\appendix

\section{Proposed Training Export}
\label{app:training-export}

Each row names an immutable information cutoff and one prospective task. The
input includes only native observations and selected world or cognitive records
visible at that cutoff, together with registry and grounding versions. The
target may be a later native observation, continuous-state interval, discrete
event, observer-perspective estimate, candidate-model prediction, representation
revision, consequence, or correction. World time, observation order,
construction or revision time, and discourse order are serialized separately
when present.

Every exported value retains its epistemic track, permitted authority-parent
paths to its unique governing context, separately identified temporal ancestry,
evidence, unit or local question, uncertainty or remainder, provenance, and
accepted/rejected status. Access and the row cutoff are applied before
traversal or serialization; denied references expose no private ancestry.
A Cut or component is exported only with the complete permitted local Cut:
parent Event, question, unit, keyed sibling answers and weights, remainder,
optional conditioning address, constraints, and recomposition rule.
If the complete view is unavailable, the Cut view is denied; an isolated
semantic weight is not a valid export. Direction weights retain their
declared proposal-policy or forecast-probability interpretation, with later
calibration recorded separately.
Model-completed intervals remain distinguishable from observations. An export
manifest reports process selection, sampling resolution, missingness, cognitive
coverage, renderer or sensor family, and any adaptive-resolution decision.
Omitted annotation means unmodeled or unannotated, not false. A snapshot-only
export is not labeled cutoff-safe until chronology and correction replay are
implemented and tested.

For joint prediction, a row aligns the next native interval $Y$, selected
process continuation $X$, discrete events $E$, and any proposed registry change
$\Delta\mathcal V$. For generative inversion, worlds, renderer families, and
grounding versions are split so that lexical or template memorization cannot
substitute for state recovery. For prospective real-data evaluation, the reveal
interval and scoring record are appended without mutating the original input or
forecast.


\subsection{Reservoir split, training, and scoring protocol}
\label{app:reservoir-protocol}

The following settings complete the protocol for
Section~\ref{sec:first-experiment}.

Generate $2{,}000$ independent world groups, each retaining its mirrored
stock/leak history, renderings, and intervention siblings. Assign
$1{,}400/300/300$ groups to train/validation/test before making rows, and use
cutoffs $8,16,32,64$. Paired variants never cross splits. Freeze a paired
ordinary-text and typed renderer for training, with identical observation
content. At test time compare the ordinary-text arms on their original
renderer and a held-out fact-preserving paraphrase renderer; keep the typed
view fixed and report that robustness comparison separately from $T-U$.
Report renderer cost and ensure no input renderer can read private targets.
Use five paired initialization seeds, batches of
$32$, AdamW at learning rate $3\times10^{-4}$, and $20{,}000$ updates per arm;
fix weight decay at $.01$, use no dropout, and report truncation (which must
be zero), token counts, parameters, and measured compute. A secondary
compute-capped comparison stops all arms at the smallest observed training
budget; equal update counts alone are not equal compute. These numbers are
design choices to freeze before test release, not a power calculation.

The primary endpoint is four-hour native forecast negative log likelihood,
averaged first within world group. Average paired $U-N$ differences across
seeds, then report a $95\%$ group-bootstrap interval and the per-seed
differences separately. Secondary outcomes include $T-U$, empty-onset Brier
score, and stock-posterior calibration. The last two apply only to $N,U,T$;
the direct native reference supplies total-reading forecasts, not a stock
posterior or tank-specific Events.
An exact oracle enumerates the $50$ initial states consistent with the visible
history and the eight future mode paths. A current-observation-only prior
baseline conditions on $(S_t,q_t)$ but discards earlier evidence. Both remain
benchmarks, never student inputs. Appendix~\ref{app:reservoir-row} gives one
row and the training/inference sequence.

\subsection{A complete reservoir forecast row}
\label{app:reservoir-row}

The following row instantiates Section~\ref{sec:first-experiment}. It is
illustrative arithmetic, not experimental evidence. Within the synthetic
world the generator knows the private targets; the student sees only the
input column. The format renderer has the same restricted access as the
student.

\begin{table}[H]
\centering
\small
\begin{tabularx}{\textwidth}{@{}p{0.19\textwidth}YY@{}}
\toprule
Field & Cutoff-visible input & Training or later scoring target \\
\midrule
Identity and clocks & World group $g$, History and two stock-process addresses;
registry digest $v$; cutoff hour $8$; horizon hours $(8,12]$ &
Same identities and source version; later reveal attached without rewriting input \\
Native history & Hours $0,\ldots,8$; totals
$(20,23,26,29,24,19,22,25,28)$ litres; modes
$(1,1,1,0,0,1,1,1,1)$; no gauge &
Future totals $(31,32,27,30)$ litres \\
Action and observation policy & Ordinary supply for all four forecast hours;
total and mode public; individual stocks and leak private &
Future modes $(q_9,q_{10},q_{11})=(1,0,1)$ \\
Selected process state & Current total $28$ litres; stocks otherwise unresolved &
$s^\star=(a_8=10,z=A)$, hence $b_8=18$; available only for the supervised loss \\
Path and Events & No future Event or contact label &
Stocks $(11,20),(12,20),(9,18),(10,20)$ at hours $9$--$12$;
$B$ reaches capacity at $9$ and $12$;
supply switches at $10$ and $11$ \\
Masks and provenance & Exact simulated observations rendered with frozen
renderer and access policy; no truncation &
$m_s=1$ for supervised arms, $0$ for native-only arm;
authored simulator truth, not observed human state; no registry revision \\
\bottomrule
\end{tabularx}
\caption{One information-separated training row. All stock numbers carry
litres; no sibling-relative semantic share is introduced.}
\label{tab:reservoir-row}
\end{table}

A minimal implementation follows four steps:
\begin{enumerate}
 \item Generate and split world groups before producing views. At each cutoff,
   assemble only the permitted observation prefix and planned action schedule;
   keep private state, future Events, and scoring data in a separate target.
 \item Encode the prefix. Normalize the joint $(a,z)$ logits over states
   compatible with the observed total and capacity bounds. The learner does
   not receive the oracle's history-filtered support or posterior.
 \item Enumerate each candidate with the eight future supply paths. Apply
   the shared decoder, group identical observable continuations, and compute
   Equation~\ref{eq:reservoir-loss}. Sum over alternatives rather than replacing
   them with independent mean stocks and a mean leak label.
 \item At inference, mask private labels entirely and save the forecast
   distribution before revealing outcomes. Score the frozen distribution and
   attach any later update as a new estimate, preserving the original.
\end{enumerate}
Other domains need their own target adapters. Measured continuous values
require a unit-aware density or declared discretization; missing targets need
masks; Cuts need complete compatible sibling sets. Nothing in this finite
example licenses adding arbitrary numerical scales to unmeasured processes.

\begingroup
\small
\sloppy
\hfuzz=1pt
\setlength{\emergencystretch}{3em}
\setlength{\multicolsep}{6pt}
\bibliographystyle{unsrt}
\bibliography{references}
\endgroup

\end{document}
